%--------------------------------------------------------------
% thesis.tex
%--------------------------------------------------------------
% Corso di Laurea in Informatica
% http://if.dsi.unifi.it/
% @Facolt\`a di Scienze Matematiche, Fisiche e Naturali
% @Universit\`a degli Studi di Firenze
%--------------------------------------------------------------
% - template for the main file of Informatica@Unifi Thesis
% - based on Classic Thesis Style Copyright (C) 2008
%   Andr\'e Miede http://www.miede.de
%--------------------------------------------------------------
\documentclass[twoside,openright,titlepage,fleqn,
	headinclude,12pt,a4paper,BCOR5mm,footinclude]{scrbook}
%--------------------------------------------------------------
\newcommand{\myItalianTitle}{Progetto e implementazione in FACPL di un monitor a runtime per il supporto al controllo continuativo degli accessi}
%\newcommand{\myItalianTitle}{Estensione del linguaggio FACPL per esprimere politiche di gestione dell'utilizzo continuativo delle risorse di un sistema di calcolo\xspace}
\newcommand{\myEnglishTitle}{Devising a run-time monitor for continuative access controls with FACPL\xspace}
% use the right myDegree option
\newcommand{\myDegree}{Corso di Laurea in Informatica\xspace}
%\newcommand{\myDegree}{
	%Corso di Laurea Specialistica in Scienze e Tecnologie
	%dell'Informazione\xspace}
\newcommand{\myName}{Filippo Mameli\xspace}
\newcommand{\myProf}{Rosario Pugliese\xspace}
\newcommand{\myOtherProf}{Andrea Margheri\xspace}
\newcommand{\mySupervisor}{\xspace}
\newcommand{\myFaculty}{Scuola di Scienze Matematiche, Fisiche e Naturali\xspace}
\newcommand{\myUni}{\protect{Universit\`a degli Studi di Firenze}\xspace}
\newcommand{\myLocation}{Firenze\xspace}
\newcommand{\myTime}{Anno Accademico 2014-2015\xspace}
\newcommand{\myVersion}{Version 0.1\xspace}
%--------------------------------------------------------------
\usepackage[italian]{babel}
\usepackage[utf8x]{inputenc}
\usepackage[T1]{fontenc}
\usepackage[square,numbers]{natbib}
\usepackage[fleqn]{amsmath}
\usepackage{ellipsis}
\usepackage{listings}
\usepackage{subfig}
\usepackage{caption}
\usepackage{appendix}
\usepackage{siunitx}
\usepackage{float}
\usepackage{blindtext}
\usepackage{scrextend}
\usepackage{enumitem}
\addtokomafont{labelinglabel}{\sffamily}
%--------------------------------------------------------------
\usepackage{pkg/dia-classicthesis-ldpkg}
%--------------------------------------------------------------
% Options for classicthesis.sty:
% tocaligned eulerchapternumbers drafting linedheaders
% listsseparated subfig nochapters beramono eulermath parts
% minionpro pdfspacing
\usepackage[eulerchapternumbers,linedheaders,subfig,beramono,eulermath,
parts]{classicthesis}
%--------------------------------------------------------------
\newlength{\abcd} % for ab..z string length calculation
% how all the floats will be aligned
\newcommand{\myfloatalign}{\centering}
\setlength{\extrarowheight}{3pt} % increase table row height
\captionsetup{format=hang,font=small}
%--------------------------------------------------------------
% Comandi personali
%--------------------------------------------------------------
\usepackage{verbatimbox}
\usepackage{listings}
\usepackage{color}
\definecolor{name}{rgb}{0.5,0.5,0.5}
\newcommand{\MyFigure}[3]{
	\begin{figure}[h]

		\centering
		\includegraphics[width=#3\linewidth,keepaspectratio]{#1}
		\caption{ #2 \label{fig:#1}}

	\end{figure}
}
\newcommand{\MyF}[4]{
	\begin{figure}[h]

		\centering
		\includegraphics[width=#3\linewidth,keepaspectratio,height=#4]{#1}
		\caption{ #2 \label{fig:#1}}

	\end{figure}
}

\newcommand{\MyFig}[4]{
	\begin{figure}[#4]

		\centering
		\includegraphics[width=#3\linewidth]{#1}
		\caption{ #2 }\label{#1}

	\end{figure}
}

\newcommand{\statusattribute}{\textit{Status Attribute }}
\newcommand{\status}{\textit{Status }}
\newcommand{\ePolicy}{\emph{Policy }}
\newcommand{\epolicy}{\emph{policy }}
\newcommand{\ePolicySet}{\emph{Policy Set }}
\newcommand{\eObligation}{\emph{Obligation }}
\newcommand{\eobligation}{\emph{obligation }}
\newcommand{\eObligations}{\emph{Obligations }}
\newcommand{\eRule}{\emph{Rule }}
\newcommand{\etarget}{\emph{target }}
\newcommand{\efalse}{\texttt{false}}
\newcommand{\etrue}{\texttt{true}}
\newcommand{\eand}{\emph{and }}
\newcommand{\eor}{\emph{or }}


\newcommand{\sr}[1]{#1}

\newcommand{\define}{\triangleq}



\newcommand{\modif}[1]{{\color{red}#1}}

%---------------------------------------
%ACRONYM
%---------------------------------------
\newcommand{\xacml}{\ac{XACML}}
\newcommand{\facpl}{\ac{FACPL}}
\newcommand{\pbac}{\ac{PBAC}}
\newcommand{\rbac}{\ac{RBAC}}

\newcommand{\pdp}{\ac{PDP}}
\newcommand{\pep}{\ac{PEP}}
\newcommand{\pap}{\ac{PAP}}
\newcommand{\pip}{\ac{PIP}}
\newcommand{\pr}{\ac{PR}}

\newcommand{\dsl}{\ac{DSL}}
\newcommand{\ide}{\ac{IDE}}

%%%%%%%%%%%%%%%%%%%%%%%%%%%%%
%SYNTAX
%%%%%%%%%%%%%%%%%%%%%%%%%%%%%
%-----------------------------------
%Algorithm
%-----------------------------------

\newcommand{\denyOver}{\x{d}\textrm{-}\x{over}_{\delta}}
\newcommand{\permitOver}{\x{p}\textrm{-}\x{over}_{\delta}}
\newcommand{\permitUnless}{\x{p}\textrm{-}\x{unless}\textrm{-}\x{d}_{\delta}}
\newcommand{\denyUnless}{\x{d}\textrm{-}\x{unless}\textrm{-}\x{p}_{\delta}}
\newcommand{\onlyOneApp}{\x{one}\textrm{-}\x{app}_{\delta}}
\newcommand{\firstApp}{\x{first}\textrm{-}\x{app}_{\delta}}
\newcommand{\weakCon}{\x{weak}\textrm{-}\x{con}_{\delta}}
\newcommand{\strongCon}{\x{strong}\textrm{-}\x{con}_{\delta}}


\newcommand{\denyOverO}[1]{\x{d}\textrm{-}\x{over}_{\x{#1}}}
\newcommand{\permitOverO}[1]{\x{p}\textrm{-}\x{over}_{\x{#1}}}
\newcommand{\permitUnlessO}[1]{\x{p}\textrm{-}\x{unless}\textrm{-}\x{d}_{\x{#1}}}
\newcommand{\denyUnlessO}[1]{\x{d}\textrm{-}\x{unless}\textrm{-}\x{p}_{\x{#1}}}
\newcommand{\onlyOneAppO}[1]{\x{one}\textrm{-}\x{app}_{\x{#1}}}
\newcommand{\firstAppO}[1]{\x{first}\textrm{-}\x{app}_{\x{#1}}}
\newcommand{\weakConO}[1]{\x{weak}\textrm{-}\x{con}_{\x{#1}}}
\newcommand{\strongConO}[1]{\x{strong}\textrm{-}\x{con}_{\x{#1}}}

%Obligation Types
\newcommand{\obM}{\modif{\x{M}}}
\newcommand{\obO}{\modif{\x{O}}}


%PEP
\newcommand{\based}{\x{base}}
\newcommand{\debugAlg}{\x{debug}}
\newcommand{\denyBiased}{\x{deny}\textrm{-}\x{biased}}
\newcommand{\permitBiased}{\x{permit}\textrm{-}\x{biased}}

%-------------------------------------
%Syntax formatting, comparison functions, examples
%-------------------------------------

\newcommand{\alice}{\x{Alice}}
\newcommand{\bob}{\x{Bob}}

\newcommand{\polSet}[4]{{\bf{\{}}#1{\,}\x{target:}\,#2 {?}\, \x{policies:}#3^{+} {\,}\x{obl:}#4^{*}\, {\bf{\}}}}
\newcommand{\polSetE}[4]{{\bf{\{}}#1 \, \x{target:}\,#2 \, \x{policies:}#3 \, \x{obl:}#4 \, {\bf{\}}}}
\newcommand{\pol}[4]{{\bf{\langle}}#1{\,}\x{target:}\,#2 {?}{\,}\, \x{rules:}#3^{+} {\,}\x{obl:}#4^{*}{\,}{\bf{\rangle}}}
\newcommand{\polE}[4]{{\bf{\langle}}#1{\,}\x{target:}\,#2 {\,}\, \x{rules:}#3 {\,}\x{obl:}#4 {\,}{\bf{\rangle}}}

\newcommand{\pdpPol}[2]{\{ #1 \,  \, #2 \}}
\newcommand{\obl}[1]{\, PepAction( #1^* ) }

\newcommand{\streq}{\x{equal}}
\newcommand{\exprOp}{\x{eop}}
\newcommand{\exprOperator}{\x{op}}

\newcommand{\attribute}[2]{( #1 , #2 )}

%\newcommand{\mEl}{m}
%\newcommand{\mEval}[2]{#2 \models #1}

\newcommand{\x}[1]{{\sf #1}}


%---------------------------------------
% Separator
%---------------------------------------

\newcommand{\set}[1]{\{#1\}} % brackets for sets
\newcommand{\Sep}{\ \mid\ }

\newcommand{\policySetBegin}{{\bf{\{}}}
\newcommand{\policyBegin}{{\bf{\{}}}
\newcommand{\policySep}{\,\bf{\ }\,}
\newcommand{\policiesBegin}{\x{policies:}\,}
\newcommand{\targetBegin}{\x{target:}\,}
\newcommand{\targetEnd}{\,{\bf{\ }}}
\newcommand{\policiesEnd}{\,{\bf{\ }}}
%\newcommand{\rulesBegin}{\x{rules:}\,}
%\newcommand{\rulesEnd}{\,{\bf{\ }}}
\newcommand{\policyEnd}{{\bf{\}}}}
\newcommand{\policySetEnd}{{\bf{\}}}}
%\newcommand{\conditionBegin}{\x{condition:}\,}
%\newcommand{\conditionEnd}{\,{\bf{\ }}}
\newcommand{\oblsBegin}{\x{obl:}\,}
\newcommand{\oblsEnd}{\bf{\ }\,}
\newcommand{\oblBegin}{{\bf{[}}}
\newcommand{\oblEnd}{{\bf{]}}}

\newcommand{\match}[3]{#1{\bf(}#2{\bf,}#3{\bf)}}
\newcommand{\ruleOpt}[1]{{\bf{(}}#1{\bf{)}}}
\newcommand{\oblOpt}[1]{{\bf{(}}#1{\bf{)}}}

\newcommand*\lfrac[2]{{}_{#1}\!\backslash\!^{#2}} % table combining algorithm

%------------------------------
%Decisions & values
%------------------------------
\newcommand{\permit}{\x{permit}}
\newcommand{\deny}{\x{deny}}
\newcommand{\notApp}{\x{not}\textrm{-}\x{app}}
\newcommand{\indet}{\x{indet}}
\newcommand{\excpt}{\perp}
\newcommand{\err}{\x{error}}

\newcommand{\indetD}{\x{indetD}}
\newcommand{\indetP}{\x{indetP}}
\newcommand{\indetDP}{\x{indetDP}}

\newcommand{\true}{\x{true}}
\newcommand{\false}{\x{false}}

%%%%%%%%%%%%%%%%%%%%%%%
%SEMANTICS
%%%%%%%%%%%%%%%%%%%%%%%

%Auxiliary Semantic Functions and notations
\newcommand{\pepSemR}[1]{(\!( #1 )\!)}
\newcommand{\concat}{{\bullet}}
\newcommand{\subobeff}[2]{#1\!\!\mid_{#2}} %sub sequence of obligations

%Name of Semantic Function
\newcommand{\denSemF}[1]{\mathcal{ #1 }}

%Semantic Brakets
%\newcommand{\denSem}[2]{[\![ #1 ]\!]_{#2}}
\newcommand{\denSem}[2]{[\![ #1 ]\!] #2}

%Named Semantics Functions
\newcommand{\policySem}[2]{\denSemF{P}\denSem{#1}{#2}}
\newcommand{\algSem}[2]{\denSemF{A}\denSem{#1}{#2}}
\newcommand{\exprSem}[2]{\denSemF{E}\denSem{#1}{#2}}
\newcommand{\oblSem}[2]{\denSemF{O}\denSem{#1}{#2}}
\newcommand{\oblSemS}[2]{\denSemF{OS}\denSem{#1}{#2}}
\newcommand{\pdpSem}[2]{\denSemF{P}dp\denSem{#1}{#2}}
\newcommand{\pepSem}[2]{\denSemF{E}A\denSem{#1}{#2}}
\newcommand{\pasSem}[1]{\denSemF{P}as\denSem{#1}{}}
\newcommand{\reqSemS}[2]{\denSemF{R}\denSem{#1}{#2}}  %two arguments, \name as second argument
\newcommand{\reqSem}[1]{\denSemF{R}\denSem{#1}{}} %used for PAS semantics

%Generic Element of each Syntactic Category
\newcommand{\expr}{\mathit{expr}}
\newcommand{\effect}{\mathit{e}}
\newcommand{\ob}{\mathit{o}}
\newcommand{\obType}{\mathit{t}}
\newcommand{\fo}{\mathit{f\!o}}
\newcommand{\foS}{\mathit{f\!o}^*}
\newcommand{\rSyntax}{\mathit{req}} %element of syntactic category Request
\newcommand{\req}{\mathit{r}} % element of functional representation of Request
\newcommand{\policy}{\mathit{p}}
\newcommand{\algSyntax}{\mathit{a}}
\newcommand{\pdpRes}{\mathit{res}}
\newcommand{\dec}{\mathit{dec}} %decision
\newcommand{\double}{\mathit{d}} %double
\newcommand{\extVal}{\mathit{w}} % Obligations and table of expression semantics
\newcommand{\enfAlg}{\mathit{ea}}
\newcommand{\pdpSyntax}{\mathit{pdp}}
\newcommand{\name}{\mathit{n}}
\newcommand{\val}{\mathit{v}}
\newcommand{\pepAction}{\mathit{pepAct}}

\newcommand{\algNT}{\mathit{Alg}} %used in syntax section
\newcommand{\algName}{\mathsf{alg}} % used in combining algorithms semantic section

%Semantics of Combiing Algorithms
\newcommand{\alg}[1]{\mathsf{alg}_{#1}}
\newcommand{\algD}{\alg{\delta}}
\newcommand{\algOp}{\otimes \mathsf{alg}}
\newcommand{\algOpAlg}[1]{\otimes #1}

\newcommand{\all}{\x{all}}
\newcommand{\greedy}{\x{greedy}}
\newcommand{\isFinal}[2]{\mathit{isFinal}_{#1}({#2})}
\newcommand{\isFinalPred}[1]{\mathit{isFinal}_{#1}}


\definecolor{eclipsePurple}{RGB}{127,0,85}
\definecolor{dkgreen}{rgb}{0,.6,0}
\definecolor{brightpink}{rgb}{1.0, 0.0, 0.5}
\definecolor{burgundy}{rgb}{0.5, 0.0, 0.13}

\renewcommand{\lstlistingname}{Codice}% Listing -> Algorithm
\renewcommand{\lstlistlistingname}{Lista di Codici}
\lstdefinelanguage{FACPL} {
  % list of keywords
	alsoletter=-,
	keywords = [1] {
		deny-biased, permit-biased, deny-unless-permit, unless, biased
	},
  	morekeywords = [2] {
		target, policies, obl, rules
	},
	keywords = [3] {
		addStatus, subStatus, equal, flagstatus,
		less-than, greater-than, sumDate, setString
	},
	keywords = [4] {
		Policy, PolicySet, Rule, Request
	},
	keywords = [5] {
		permit, deny
	},
  %style
	keywordstyle = [1]{
		\color{dkgreen}
	},
	keywordstyle = [2]{
		\color{burgundy}\bfseries
	},
	keywordstyle = [3]{
		\color{dkgreen}
	},
	keywordstyle = [4]{
		\color{orange}\bfseries
	},
	keywordstyle = [5]{
		\color{brightpink}\bfseries
	},
	numbers=left,
	sensitive=false,
	morecomment=[l]{//},
	morecomment=[s]{/*}{*/},
	morestring=[b]",
	morestring=[b]',
	stringstyle=\color{blue}\ttfamily
}
%--------------------------------------------------------------
% Java
%--------------------------------------------------------------
\definecolor{mygray}{rgb}{0.5,0.5,0.5}


\definecolor{monokai@black}{HTML}{2C2C2A}
\definecolor{monokai@gray}{HTML}{524F52}
\definecolor{monokai@yellow}{HTML}{FDF5A9}
\definecolor{monokai@orange}{HTML}{FD971F}
\definecolor{monokai@red}{HTML}{DC322F}
\definecolor{monokai@magenta}{HTML}{D33682}
\definecolor{monokai@violet}{HTML}{F7208B}
\definecolor{monokai@blue}{HTML}{004C99}
\definecolor{monokai@green}{HTML}{009900}

\lstdefinestyle{monokai}{
	frame=tb,
	language=Java,
	aboveskip=3mm,
	belowskip=3mm,
	showstringspaces=false,
	columns=flexible,
	basicstyle={\small\ttfamily},
	numbers=left,
	numberstyle=\tiny\color{monokai@black},
	numbersep=5pt,
	stepnumber=2,
	breaklines=true,
	breakatwhitespace=true,
	tabsize=3,
	numberstyle=\color{monokai@gray},
	keywordstyle=\color{monokai@violet},
	stringstyle=\color{monokai@green}\ttfamily,
	identifierstyle=\color{monokai@black},
	commentstyle=\color{monokai@blue},
	emphstyle=\color{monokai@red}
}

\lstdefinestyle{javigno}{
	frame=tb,
	language=Java,
	aboveskip=3mm,
	belowskip=3mm,
	showstringspaces=false,
	columns=flexible,
	basicstyle={\small\ttfamily},
	numbers=left,
	numberstyle=\tiny\color{mygray},
	numbersep=5pt,
	stepnumber=2,
	breaklines=true,
	breakatwhitespace=true,
	tabsize=3
}

\newcommand{\javaFirstline}[4]
	{
		\lstinputlisting[language = Java, style = monokai, firstline = #2, lastline = #3,firstnumber=1]{./code/cap5/#1}\label{lst:#4}
	}
\newcommand{\javaLastline}[4]
	{
		\lstinputlisting[language = Java, style = monokai, firstline = #2, lastline = #3,firstnumber=last]{./code/cap5/#1}\label{lst:#4}
	}
\newcommand{\javaFirstlineN}[5]
	{
		\lstinputlisting[language = Java, style = monokai, firstline = #2, lastline = #3,firstnumber=#5]{./code/cap5/#1}\label{lst:#4}
	}
\newcommand{\javaFirstlineC}[5]
	{
		\lstinputlisting[language = Java, style = monokai, firstline = #3, lastline = #4,firstnumber=1, caption = {#2}\label{lst:#5}]{./code/cap5/#1}
	}
\newcommand{\javaLastlineC}[5]
	{
		\lstinputlisting[language = Java, style = monokai, firstline = #3, lastline = #4,firstnumber=last, caption = {#2}\label{lst:#5}]{./code/cap5/#1}
	}
\newcommand{\javaFirstlineNC}[6]
	{
		\lstinputlisting[language = Java, style = monokai, firstline = #3, lastline = #4,firstnumber=#6, caption = {#2}\label{lst:#5}]{./code/cap5/#1}
	}
%--------------------------------------------------------------
% Layout setting
%--------------------------------------------------------------
\usepackage{geometry}
\geometry{
	a4paper,
	ignoremp,
	bindingoffset = 1cm,
	textwidth     = 13.5cm,
	textheight    = 21.5cm,
	lmargin       = 3.5cm, % left margin
	tmargin       = 4cm    % top margin
}

\lstset{
  	frame=tb,
		language=Matlab,
  	aboveskip=3mm,
  	belowskip=3mm,
  	showstringspaces=false,
  	columns=flexible,
  	basicstyle={\small\ttfamily},
  	numbers=none,
  	breaklines=true,
  	breakatwhitespace=true,
  	tabsize=3
}
%--------------------------------------------------------------
\begin{document}
\frenchspacing
\raggedbottom
\pagenumbering{roman}
\pagestyle{plain}
%--------------------------------------------------------------
% Frontmatter
%--------------------------------------------------------------
%--------------------------------------------------------------
% titlepage.tex (use thesis.tex as main file)
%--------------------------------------------------------------
\begin{titlepage}
	\begin{center}
   	\large
      \hfill
      \vfill
      \begingroup
         \includegraphics[scale=0.15]{logo/LOGO}\\
%			\spacedallcaps{\myUni} \\
			\myFaculty \\
			\myDegree \\
			\vspace{0.5cm}
         \vspace{0.5cm}
         Tesi di Laurea
      \endgroup
      \vfill
      \begingroup
      	\color{Maroon}\spacedallcaps{\myItalianTitle} \\ $\ $\\
      	\spacedallcaps{\myEnglishTitle} \\
	\bigskip
      \endgroup
      \spacedlowsmallcaps{\myName}
      \vfill
      \vfill
      Relatore: \emph{Rosario Pugliese}\\
      Correlatore: \emph{Andrea Margheri}\\
      \vfill
      \vfill
      \myTime
      \vfill
	\end{center}
\end{titlepage}
%--------------------------------------------------------------
% back titlepage
%--------------------------------------------------------------
   \newpage
	\thispagestyle{empty}
	\hfill
	\vfill
	\noindent\myName:
	\textit{\myItalianTitle,}
	\myDegree, \textcopyright\ \myTime
%--------------------------------------------------------------
% back titlepage end
%--------------------------------------------------------------

\pagestyle{scrheadings}
%--------------------------------------------------------------
% Mainmatter
%--------------------------------------------------------------
\pagenumbering{arabic}
% use \cleardoublepage here to avoid problems with pdfbookmark
%\include{intro} % use \myChapter command instead of \chapter
\tableofcontents
\listoffigures
\cleardoublepage
%\thispagestyle{empty}
%\begin{flushright}
%\null\vspace{\stretch {1}}
%\emph{"Inserire citazione" \break --- Inserire autore citazione} \vspace{\stretch{2}}\null
%\end{flushright}
%\cleardoublepage
\myChapter{Introduzione}
I sistemi informatici, fin dalla loro nascita, sono stati utilizzati per la gestione di dati.
Il tipo di informazioni gestite ha reso necessario la creazione di un sistema che le proteggesse, infatti i
dati più sensibili, se diffusi senza una valida autorizzazione, possono arrecare ad esempio danni economici
ad una società o nuocere anche gli utenti nel privato.
Per far fronte al problema sono stati sviluppati dei modelli per il controllo degli accessi. Tuttavia la quantità di dati
e la complessità dei sistemi moderni ha mostrato i limiti delle tecnologie concepite e questo fatto ha portato inevitabilmente
allo sviluppo di un nuovo approccio: lo \ac{UCON}.

I ricercatori dell'università di San Antonio, Jaehong Park e Ravi Sandhu, hanno descritto un modello
di \ac{UCON}\cite{ucon} che consente di ottenere decisioni durante l'accesso e di basare le valutazioni di una richiesta
sul comportamento passato.

In questa tesi si descrive un'estensione basata su UCON per \ac{FACPL}.
Il problema principale da risolvere è la staticità del processo di valutazione del linguaggio.
Tutte le richieste di accesso ricevono una risposta attraverso un sistema di verifica sempre uguale e i controlli
risultano ridondati e in alcuni casi inutili.
Nella prima parte dello sviluppo è stata modificata la sintassi del linguaggio per adeguarla ad una nuova
gestione delle richieste, in seguito
è stata sviluppata un'estensione della libreria su cui si basa \ac{FACPL} per modificare
il processo valutativo al seguito di richieste di accesso continuative.
Nei capitoli successivi si mostrano i passi che hanno portato allo sviluppo della nuova gestione delle richieste
in grado di rendere dinamico il processo di valutazione.
\vspace{5mm}
\newline
Il resto del documento è così strutturato:
\begin{itemize}
  \renewcommand\labelitemi{--}
  \item Nel Capitolo~\ref{chap:Access Control e Usage Control}
  si introducono vari modelli di Access Control, si descrive lo \ac{UCON} e si mostrano due casi di studio basati sul nuovo modello.
  \item Nel Capitolo~\ref{chap:Formal Access Control Policy Language}
  si descrive la sintassi e la semantica di \ac{FACPL}, il processo di valutazione e un esempio utilizzando il linguaggio.
  \item Nel Capitolo~\ref{chap:Usage Control in FACPL}
  si mostra l'estensione di \ac{FACPL} per migliorare la gestione delle richieste.
  \item Nel Capitolo~\ref{chap:Estensione della libreria FACPL}
  si riportano le modifiche alla libreria Java su cui è basato il linguaggio per adattarlo al nuovo processo valutativo.
  \item Nel Capitolo~\ref{chap:Conclusioni}
  si riassume il lavoro svolto e si presentano alcuni possibili sviluppi futuri.
\end{itemize}

\myChapter{Access Control e Usage Control}
\label{chap:Access Control e Usage Control}
L'esigenza di proteggere i dati ha determinato la necessità di creare strumenti per il controllo degli accessi così da eliminare,
 o almeno limitare, i rischi derivati dall'accesso non autorizzato alle informazioni.\par
Nel corso del tempo si sono sviluppati numerosi modelli per i sistemi di controllo degli accessi~\cite{NISTACM}.
Nelle sezioni successive saranno presentati:
\ac{ACL}, \ac{RBAC}, \ac{ABAC} e \ac{PBAC}. Questi modelli sono introdotti in Sezione~\ref{sub:PBAC}.
In Sezione~\ref{sec:Usage_Control} è descritto lo \ac{UCON} e il recente modello UCON\textsubscript{ABC}, poi in Sezione~\ref{subs:es_UC}
sono esposti due esempi che utilizzano \ac{UCON}.
\section{Access Control}

\subsection*{Access Control List}
\label{sub:ACL}
\ac{ACL} è stato creato agli inizi degli anni settanta per la necessità di aver un controllo degli
accessi sui sistemi multiutente.
Utilizza una lista di utenti con annesse le possibili azioni autorizzate~\cite{ACL}. Il modello è semplice,
ma ha molte limitazioni. Quando nel sistema ci sono numerosi utenti o risorse, la quantità di dati da verificare diventa
difficile da gestire. Questo può portare a errori di assegnazione di autorizzazioni e ad un eccessivo numero di controlli
necessari per un singolo accesso.
\subsection*{Role Based Access Control}
\label{sub:RBAC}
\ac{RBAC} è l'evoluzione di \ac{ACL}. In questo modello vengono introdotti i \emph{ruoli}~\cite{RBAC}.
Il ruolo è un titolo che definisce un livello di autorità. Più utenti
possono avere lo stesso ruolo e quindi avere a disposizione tutte le risorse connesse a questo. Il modello
diventa scalabile e più facile da gestire, inoltre si possono anche creare delle gerarchie per facilitare l'assegnamento di
risorse in base alla classificazione dell'utente. \ac{RBAC} è limitato dalla staticità dei permessi e dal concetto di ruolo
che non può essere associato a nessun attributo.
\par In Figura \ref{fig:img/RBAC} si mostra un gruppo di utenti
facente parte di un \emph{role} e la connessione tra i permessi e il ruolo.
\MyFigure{img/RBAC}{RBAC}{0.7}
\subsection*{Attribute Based Access Control}
\label{sub:ABAC}
\ac{ABAC} si basa sull'utilizzo di attributi associati all'utente, all'azione o al contesto
della richiesta~\cite{ABAC}. La valutazione di una autorizzazione diventa più specifica rispetto a \ac{RBAC},
le regole sono più precise per ogni risorsa e inoltre il modello si basa su politiche invece che su permessi statici.

In Figura \ref{fig:img/abac.png} si illustra l'utilizzo da parte del sistema di valutazione~(Rules engine) di
attributi dell'\emph{utente}~(user), del \emph{contesto}~(environment) e delle \emph{azioni}, per la decisione
sull'autorizzazione.
\MyFigure{img/abac.png}{ABAC}{0.7}

Il modello non è utilizzato nei sistemi operativi, dove \ac{ACL} e \ac{RBAC} sono i più diffusi, ma è sviluppatto
spesso a livello applicativo. Il problema fondamentale di questo paradigma è che le regole non sono uniformi e se il numero
di risorse è consistente, la gestione di queste diventa complicata. Il modello \acf{PBAC} cerca di
risolvere i difetti di \ac{ABAC}.
\subsection*{Policy Based Access Control}
\label{sub:PBAC}
\ac{PBAC} riorganizza il modello \ac{ABAC} per semplificare la gestione delle regole~\cite{NISTACM}.
\MakeUppercase{è} una standardizzazione di \ac{ABAC} e si basa su \emph{politiche} che non sono altro che insiemi di \emph{rule}.
\ac{PBAC} combina gli attributi della risorsa, dell'ambiente e dell'utente, e li utilizza insieme alle \emph{rule} per
determinare l'autorizzazione.
Le politiche possono essere messe insieme per creare gruppi di politiche, in questo modo il sistema diventa scalabile e di
più facile utilizzo.\par
Per costruire un sistema di controllo degli accessi basato sul modello \ac{PBAC} è necessario però l'utilizzo di
un linguaggio adatto allo scopo. \ac{OASIS} ha creato il linguaggio \ac{XACML} che è diventato lo standard per lo sviluppo di un sistema costruito sul modello PBAC.
\section{Usage Control}
\label{sec:Usage_Control}
Dopo quaranta anni di studi sul controllo degli accessi i modelli sviluppati si sono consolidati e sono largamente utilizzati
su sistemi operativi o applicazioni. Tuttavia la complessità e la varietà degli ambienti informatici moderni va oltre i limiti
dei modelli creati ~\cite{uconSN}.\par
Il termine Usage Control \ac{UCON} è stato ripreso da Jaehong Park e Ravi Sandhu per creare il modello UCON\textsubscript{ABC}~\cite{ucon}.
Questo è una generalizzazione dell'Access Control che include obbligazioni, condizioni sull'utilizzo,
controlli continuativi e mutabilità ~\cite{UCRT}. Comprende e migliora i modelli di controllo di accesso tradizionali aggiungendo
obblighi da parte dell'utente per accedere ad una risorsa, vincoli sulle azioni permesse,
la gestione di attributi variabili nel tempo e la continuità nella valutazione delle decisioni.\par
Il modello UCON\textsubscript{ABC} estende i controlli sull'accesso tradizionali ed è composto da otto componenti fondamentali.
Queste sono subject, subject attribute, object, object attribute, right, authorization, obligation e condition.\par
\MyFigure{img/ucon.png}{UCON}{1.1}
I \emph{Subject} sono entità a cui si associano degli attributi e hanno o esercitano \emph{Right} sugli \emph{Object}. Possiamo per semplicità
associare i Subject ad un singolo individuo umano.\par
Gli \emph{Object} sono insiemi di entità su cui i \emph{Subject} possono avere dei \emph{Right}, questi possono essere usati o vi si può
fare accesso. Possono essere associati ad esempio a un libro, o a una qualsiasi risorsa.\par
I \emph{Right} sono i privilegi che i \emph{Subject} hanno o esercitano sugli \emph{Object}.\par
I tre fattori Authorization, oBligation e Condition (da cui prende anche il
nome il modello UCON\textsubscript{ABC}) sono predicati funzionali che devono essere valutati per le decisioni sull'uso.
I tradizionali modelli di Access Control utilizzano solo le Authorization per il processo di decisione, Obligation e Condition
sono i nuovi componenti che entrano a far parte della valutazione ~\cite{IEEE}.\par
Le \emph{Authorization} devono valutare la decisione sull'uso. Queste danno un responso positivo
o negativo a seconda che la domanda di un Subject sia accettata o meno.\par
Le \emph{Obligation} verificano i requisiti obbligatori che un Subject deve eseguire prima o durante l'utilizzo di una risorsa.\par
Infine le \emph{Condition} restituiscono \texttt{true} o \texttt{false} in base alle variabili dell'ambiente o allo stato del sistema.\par
Il processo di decisione è diviso in tre fasi come mostrato in Figura \ref{fig:img/uconFase.png} : Before usage (pre), Ongoing usage (on) e After usage ~\cite{SurveyUsageControl}.
\MyFigure{img/uconFase.png}{Fasi del processo di decisione}{1}\\
La valutazione della prima parte inizia da una richiesta e non ha differenze con il processo valutativo dell'Access Control.
Nella seconda invece si utilizzano i nuovi predicati introdotti ed è in questa parte che si affermano i controlli continuativi,
le obbligazioni e le condizioni sull'utilizzo.\par
L'ultima fase varia in base agli eventi delle fasi precedenti. Ad esempio se il Subject che ha richiesto un accesso ha violato
una policy oltre al non aver ricevuto l'autorizzazione potrebbe anche essere ammonito e il sistema potrebbe non accettare più
nessuna sua richiesta.

\section{Esempi di Usage Control}
\label{subs:es_UC}
Si propongono adesso due esempi in cui si utilizza lo Usage Control
\subsubsection*{Accessi in lettura e scrittura di file}
\label{subs:Accessi in lettura e scrittura di file}
In un sistema ci sono alcuni file e gli utenti sono divisi in gruppi, si suppone inoltre che le richieste di lettura
siano più comuni rispetto a quelle di scrittura.\par Tutti possono leggere i file, ma solo
gli utenti nel gruppo degli amministratori possono modificare un elemento. Se uno degli amministratori sta scrivendo
su un file, nessuno può leggere o apportare modifiche nello stesso momento. Una volta finita la scrittura si potrà sbloccare
i file e renderli di nuovo disponibili a tutti.\par
Letture ripetute, su uno stesso insieme di file, vengono gestite in modo efficiente.
\subsubsection*{Servizio di streaming}
\label{subs:Servizio di streaming}
In un sistema gli utenti, dopo aver fatto il login, possono fare una richiesta di ascolto di un brano. I clienti
sono divisi in due tipologie. Chi effettua il login come utente premium può fare richieste di ascolto senza nessuna restrizione.
I clienti stardard hanno invece un tempo limite esaurito il quale devono ascoltare la pubblicità prima di poter fare
un'altra richiesta di ascolto. In questo caso si assumono le richieste di ascolto più numerose rispetto ai login o ai logout.\par
\vspace{1em}
I due esempi non possono essere implementati con i modelli di Access Control tradizionali perché basano
le valutazioni sul comportamento passato e richieste ripetute di uno stesso tipo sono valutate
in modo differente.

\myChapter{Formal Access Control Policy Language }
\label{chap:Formal Access Control Policy Language}
Il linguaggio \ac{FACPL} (fakpol) è stato creato come alternativa a \ac{XACML}. Come accennato nel in \ref{sub:PBAC},
l'organizzazione \ac{OASIS} ha ideato il linguaggio \ac{XACML} per sviluppare sistemi basati sul modello \ac{PBAC}.\par
\ac{XML} è utilizzato da \ac{XACML} per definire le sue politiche. Questo linguaggio di markup è molto usato per lo scambio
di dati tra sistemi, ma rende le politiche difficili da comprendere. Infatti anche le regole più semplici sono prolisse
e questo rende la lettura problematica per un utente.\par
\ac{FACPL} nasce dalle idee di \ac{XACML} e ha l'obiettivo di semplificare la scrittura di politiche e di definire un framework
costruito sopra basi formali solide, in modo da permettere agli sviluppatori di specificare e verificare automaticamente
delle proprietà.

In questo capitolo si delinea prima di tutto il processo di valutazione in Sezione \ref{sec:Il processo di valutazione Facpl}
poi si mostrano le componenti del sistema in Sezione \ref{sub:Componenti del sistema} e il loro significato in Sezione \ref{sec:sem_fpl}. Infine si
mostrerà in Sezione \ref{sec:Esempio_FACPL} con un esempio l'utilizzo di \ac{FACPL}.
\section{Il processo di valutazione}
\label{sec:Il processo di valutazione Facpl}
In questa sezione si espone la sequenza di azioni che il sistema compie affinchè una richiesta
di accesso ad una risorsa in input sia valutata.
\MyFigure{img/evaluationProcess.png}{Processo di valutazione}{1}
In Figura \ref{fig:img/evaluationProcess.png} si mostrano i passi che vengono eseguiti.
Le componenti chiave del processo sono tre:
\begin{itemize}
  \renewcommand\labelitemi{--}
  \item \ac{PR}
  \item \acf{PDP}
  \item \acf{PEP}
\end{itemize}
Si assume che ad ogni risorsa siano associate una o più politiche, e che queste definiscano le credenziali
necessarie per ottenere l'accesso. Il \ac{PR} contiene le Policy e le rende disponibili al \ac{PDP} (Step 1), il
quale decide se l'accesso viene autorizzato o meno. \par
Quando la richiesta è ricevuta dal \ac{PEP} (Step 2), le sue credenziali vengono codificate in una sequenza
di attributi (una coppia nome-valore) per poi utilizzare quest'ultima per creare la richiesta \ac{FACPL} (Step 3).\par
Il Context Handler aggiunge attributi dell'ambiente (come la temperatura esterna) alla richiesta e
poi invierà questa al \ac{PDP} (Step 4).\par
Il processo di autorizzazione del \ac{PDP} restituisce una \emph{risposta} verificando che gli attributi, che possono far
parte della richiesta o del contesto, siano conformi ai controlli delle policy (Step 5-8). La \emph{risposta} del \ac{PDP}
contiene una \emph{decisione} e una possibile \emph{obbligazione} (Step 9-10).\par
La \emph{decisione} può essere di quattro tipi:
\begin{itemize}
  \renewcommand\labelitemi{--}
  \item Permit
  \item Deny
  \item Not-applicable
  \item Indeterminate
\end{itemize}
I primi due tipi indicano rispettivamente richiesta accettata e richiesta non accettata. Se viene restituito Not-applicable
non ci sono policy su cui si poteva valutare la decisione, se ci sono altri errori la risposta è Indeterminate.
Le \ePolicy gestiscono automaticamente il risultato Indeterminate combinandolo con le altre decisioni
a seconda della strategia del \emph{Combining Algorithm} associato alla Policy stessa.\par
Le \emph{obligation} sono azioni addizionali che devono essere eseguite dopo la restituzione di una decisione.
Solitamente corrispondono all'aggiornamento di un file, l'invio di un messaggio o l'esecuzione di un comando.
Il \ac{PEP} ha compito di controllare l'esecuzione delle \emph{Obligation} tramite l'\emph{obligation service} (Step 11-12).
Il processo di \emph{enforcement} eseguito dal \ac{PEP} determina l'\emph{enforced decision} in base al risultato delle obbligazioni.
Questa decisione può differire da quella del \ac{PDP} e corrisponde alla valutazione finale restituita dal processo.


\section{Sintassi}
\label{sub:Componenti del sistema}
Nella Tabella \ref{tab:facpl_syntax} si mostra la sintassi di \ac{FACPL}. Questa è data da una grammatica di tipo \ac{EBNF} dove il simbolo ? indica
elementi opzionali, * sequenze (anche vuote) e + sequenze non vuote.\par
\begin{table}[h]
\centering
\scriptsize
\caption{Sintassi di FACPL}
\hrule
$
\begin{array}{@{\,}r@{\ \ }r@{\ }r@{\ \ }l@{\ }}
{\textbf{Policy Authorisation Systems}} &
\mathit{PAS} & ::= & ( \,  \x{pep:} \, \mathit{EnfAlg}\ \ \x{pdp:}\, \mathit{PDP} \, )
\\[.2cm]
{\textbf{Enforcement algorithms}} &
\mathit{EnfAlg}
& ::= & \based \Sep \denyBiased \Sep \permitBiased
\\[.2cm]
{\textbf{Policy Decision Points}} &
\mathit{PDP} & ::= & \pdpPol{\algNT\ }{\x{policies:} \, \mathit{Policy}^{+}}
\\[.2cm]
{\textbf{Combining algorithms}} &
\algNT & ::= & \permitOver \Sep \denyOver \Sep \denyUnless \Sep \permitUnless \\
&& \mid &
\firstApp \Sep \onlyOneApp \Sep \weakCon \Sep \strongCon
\\[.2cm]
{\textbf{Policies}} &
\mathit{Policy} & ::= &
\ruleOpt{\mathit{Effect}\ \ \x{target:} \, Expr\ \ \x{obl:} \, \mathit{Obligation}^{*} \, } \\
&& \mid &
\{ \algNT\ \ \x{target:} \, Expr\ \
\x{policies:} \, \mathit{Policy}^{+} \ \ \x{obl:} \, \mathit{Obligation}^{*} \, \}
\\[.2cm]
{\textbf{Effects}} &
\mathit{Effect} & ::= & \permit \Sep \deny
\\[.2cm]
{\textbf{Obligations}} &
\mathit{Obligation} & ::= & [ \, \mathit{Effect} \ \ \mathit{ObType} \ \ \obl{Expr} \, ]
\\[.2cm]
{\textbf{Obligation Types}} &
\mathit{ObType} & ::= & M \Sep O
\\[.4cm]
\textbf{Expressions}&
\mathit{Expr} & ::= &
\mathit{Name} \Sep \mathit{Value}  \\
& & \mid &\x{and(\mathit{Expr}, \mathit{Expr})} \Sep \x{or(\mathit{Expr}, \mathit{Expr})} \Sep \x{not(\mathit{Expr})} \\
& & \mid &
 \x{equal(\mathit{Expr},\mathit{Expr})}  \Sep \x{in}(\mathit{Expr}, \mathit{Expr}) \\
& & \mid & \x{greater}\textrm{-}\x{than(\mathit{Expr},\mathit{Expr})} \Sep \x{add(\mathit{Expr} ,\mathit{Expr} )}\\
& & \mid & \x{subtract(\mathit{Expr} ,\mathit{Expr} )} \Sep \x{divide(\mathit{Expr} ,\mathit{Expr} )}\\
& & \mid & \x{multiply(\mathit{Expr} ,\mathit{Expr} )} \\[.2cm]
%
\textbf{Attribute Names} &
\mathit{Name} & ::= & \mathit{Identifier}/\mathit{Identifier} \\[.2cm]
%
\textbf{Literal Values} &
\mathit{Value} & ::= & \x{true} \mid \x{false} \mid \mathit{Double} \mid \mathit{String} \mid \mathit{Date}
\\[.4cm]
{\textbf{Requests}} &
\mathit{Request} & ::= & {\attribute{\mathit{Name}}{\mathit{Value}}}^{+}
\\[.1cm]
\end{array}
$
\hrule
\label{tab:facpl_syntax}
\end{table}
{\label{Sintassi di FACPL}}
Al livello più alto troviamo il termine  \ac{PAS} che comprende le specifiche del \ac{PEP} e del \ac{PDP}.\par
Il \ac{PEP} è definito con un \emph{enforcing algorithm} che sarà applicato per stabilire quali sono le decisioni che devono
passare al processo di \emph{enforcement}.\par
Il \ac{PDP} è definito come una sequenza di \emph{Policy}\textsuperscript{+} e da un algoritmo \emph{Alg} per combinare i
risultati della valutazione delle policy.\par
Una \ePolicy può essere una \eRule semplice oppure un \ePolicySet cioè un insieme di rule o altri
policy set. In questo modo si possono creare regole singole, ma anche gerarchie di regole.\par
Un \ePolicySet è definito da un \etarget che indica l'insieme di richieste di accesso al quale la policy
viene applicata, da una lista di \emph{obligation} cioè le azioni opzionali o obbligatoriche da eseguire, da una
sequenza di \ePolicy e da un algoritmo per la combinazione dei risultati delle politiche da cui è composto l'insieme.\par
Una \eRule  è specificata da un \emph{effect} , che può essere permit o deny , da un \etarget
e da una lista di \emph{obligation}(che può essere anche vuota). \par
Le \emph{Expressions} sono costituite da \emph{attribute names} e da valori letterali(per esempio booleani, stringhe, date).\par
Un \emph{attribute name} indica il valore di un attributo. Questo può essere contenuto in una richiesta o nel contesto. La
struttura di un attribute name è della forma \emph{Identifier}/\emph{Identifier}. Dove il primo elemento indica la categoria
e il secondo il nome dell'attributo. Per esempio Name / ID rappresenta il valore di un attributo ID di categoria Name.\par
Un \emph{combining algorithm} ha lo scopo di risolvere conflitti delle decisioni date dalle valutazioni delle policy.
Ad esempio permit-overrides assegna la precedenza alle decisioni con effetto permit rispetto a quelle di tipo deny.\par
Una \emph{obligation} è definita da un effetto, da un tipo (M per obbligatorio e O per opzionale) e da un'azione e le relative
espressioni come argomento.\par

Una richiesta consiste in una sequenza di \emph{attribute} organizzati in categorie.\par
La risposta ad una richiesta \ac{FACPL} è scritta utilizzando la sintassi ausiliaria riportata in Tabella \ref{tab:facpl_context_syntax}.
La valutazione in due passi descritta in \ref{sec:Il processo di valutazione Facpl} produce due tipi differenti di risposte:
\begin{itemize}
  \renewcommand\labelitemi{--}
  \item \emph{\ac{PDP} Response}
  \item \emph{Decisions}
\end{itemize}
La prima, nel caso in cui la decisione sia permit o deny, si associa a una sequenza (anche vuota) di fulfilled obligation.\par
Una \emph{Fulfilled Obligation} è una coppia formata da un tipo e da un'azione con i rispettivi argomenti risultato della
valutazione del \ac{PDP}.
\begin{table}[!t]
\small
\caption{Sintassi ausiliaria per le risposte}
$
\begin{array}{@{\ }r@{\ \ \ \ }r@{\ }r@{\ \ }l@{\ }}

&&&\\[-.2cm]
{\textbf{PDP \ Responses }} &
\mathit{PDPResponse} & ::= & \langle \,\mathit{Decision} \ \ \ \mathit{FObligation}^* \rangle
\\[.2cm]
{\textbf{Decisions}} &
\mathit{Decision} & ::= & \permit \Sep \deny \Sep \notApp \Sep \indet
\\[.2cm]
{\textbf{Fulfilled obligations}} &
\mathit{FObligation} & ::= &  [ \, \mathit{ObType} \ \ \obl{\mathit{Value}} \, ]\\[.1cm]

\end{array}
$
\label{tab:facpl_context_syntax}
\end{table}


\section{Semantica}
\label{sec:sem_fpl}
Si presenta adesso informalmente il processo di autorizzazione del \ac{PDP} e poi quello di enforcement del \ac{PEP}.\par
Quando il \ac{PDP} riceve una richiesta di accesso, valuta la richiesta in base alle \epolicy disponibili poi determina
il risultato unendo le decisioni restituite dalle policy con l'utilizzo dei \emph{combining algorithms}.\par
La valutazione di una policy inizia dalla verifica dell'applicabilità della richiesta, che è compiuta valutando
l'espressione definita dal \emph{target}. Ci sono due casi:
\begin{description}
  \item[\normalfont\bfseries{\MakeUppercase{L}a verifica dà un risultato positivo.}] \hfill \\
        Nel caso in cui ci siano \emph{rules}, l'effetto della regola viene restituito.
        Nel caso di \ePolicySet, il risultato è ottenuto dalla valutazione delle \epolicy contenute e dalla combinazione di queste,
        tramite l'algoritmo specificato dal \ac{PDP}. In entrambi i casi in seguito si procederà al \emph{fulfilment} delle
        \emph{obligation} da parte del \ac{PEP}.
  \item[\normalfont\bfseries{\MakeUppercase{L}a verifica dà un risultato negativo.}] \hfill \\
        Viene restituito \texttt{not-app}. Nel caso di \texttt{ERROR} o di un valore non booleano,
        si restituisce \texttt{indet}.
\end{description}
Valutare le espressioni corrisponde ad applicare gli operatori, a determinare le occorrenze degli \emph{attribute names} e
a ricavarne il valore associato.\par
Se questo non è possibile, ad esempio l'attibuto è mancante e non può essere recuperato dal \emph{context handler}, si
restituisce il valore speciale $\bot$. Questo valore è usato per implementare strategie diverse per gestire
la mancanza di un attributo. $\bot$ è trattato da \ac{FACPL} in modo simile a \texttt{false}, si gestiscono
così gli attributi mancanti senza generare errori.\par
Gli operatori tengono conto degli argomenti e dei valori speciali come $\bot$ e \texttt{ERROR}.
Se gli argomenti sono \texttt{true} oppure \texttt{false} gli operatori sono applicati in modo regolare.
Se c'è un argomento $\bot$ e non ci sono \texttt{ERROR} si restituisce $\bot$, \texttt{ERROR} altrimenti.
Gli operatori \eand e \eor trattano diversamente i valori speciali. Specificatamente, \eand restituisce \texttt{true} se entrambi
gli operandi sono \etrue, \texttt{false} se almeno uno è \efalse, $\bot$ se almeno uno è $\bot$ e nessun altro è
\texttt{false} o \texttt{ERROR}, \texttt{ERROR} negli altri casi. L'operatore \eor è il duale di \eand e ha priorità minore.
L'operatore unario \emph{not} cambia i valori di \texttt{true} e \texttt{false}, ma non quelli di $\bot$ e \texttt{ERROR}.\par
La valutazione di una \epolicy termina con il \emph{fulfilment} di tutte le \emph{obligation}. Questo consiste nel valutare
tutti gli argomenti dell'azione corrispondente all'\emph{obligation}. All'occorrenza di un errore, la decisione della \epolicy
sarà modificata in \texttt{indet}.
Negli altri casi la decisione non cambierà e sarà quella del \ac{PDP} prima del \emph{fulfillment}.\par
Per valutare gli insiemi di \epolicy si devono applicare dei \emph{combining algorithm} specifici. Data una sequenza di
policy in input gli algoritmi stabiliscono una sequenza di valutazioni per le \epolicy date. Si propone l'algoritmo
\emph{permit-overrides} come esempio.
\begin{description}[labelindent=5pt,style=multiline,leftmargin=3cm]
\item[permit-overrides]
                        Se la valutazione di una \epolicy restituisce \emph{permit}, allora il risultato è \emph{permit}.
                        In altre parole, \emph{permit} ha la precedenza, indipendente dal risultato delle altre \emph{policy}.
                        Invece, se c'è almeno una \epolicy che restituisce \emph{deny} e tutte le altre restituiscono
                        \texttt{not-app} o \emph{deny}, allora il risultato è \emph{deny}. Se tutte le policy restituiscono
                        \texttt{not-app}, allora il risultato è \texttt{not-app}. In tutti gli altri casi, il risultato è \texttt{indet}.
\end{description}
Se il risultato della decisione è \emph{permit} o \emph{deny}, ogni algoritmo restituisce una sequenza di \emph{fulfilled obbligations}
conforme alla strategia di \emph{fulfilment} scelta, in seguito indicata con $\delta$. Ci sono due possibili strategie:
\begin{description}[labelindent=5pt,style=multiline,leftmargin=3cm]
  \item[All]La strategia \emph{all} richiede la valutazione di tutte le \epolicy appartenenti alla sequenza in input e
            restituisce le \emph{fulfilled obligation} relative a tutte le decisioni.
  \item[Greedy]La strategia \emph{greedy} stabilisce che, se ad un certo punto, la valutazione della sequenza di \epolicy
               in input non può più cambiare, si può non esaminare le altre policy e terminare l'esecuzione. In questo modo
               si migliorano le prestazioni della valutazione in quanto non si sprecano risorse computazionali
               per valutazioni di policy che non avrebbero alcun impatto sul risultato finale.
\end{description}\par
L'ultimo passo consiste nell'inviare la risposta del \ac{PDP} al \ac{PEP} per l'\emph{enforcement}. A questo scopo, il \ac{PEP} verifica
che tutti gli eventuali obblighi imposti dal \ac{PDP} al richiedente siano soddisfatti e decide, in base all'algoritmo di
enforcement scelto, il comportamento per le decisioni di tipo \texttt{not-app} e \texttt{indet}. Gli algoritmi sono:
\begin{description}[labelindent=5pt,style=multiline,leftmargin=4cm]
  \item[base]Il \ac{PEP} mantiene tutte le decisioni, ma se c'è un errore nella verifica degli obblighi il risultato è \texttt{indet}.
  \item[deny-biased]Il \ac{PEP} concede l'accesso solo nel caso in cui questa sia la decisione del \ac{PDP} e gli eventuali obblighi imposti dal \ac{PDP} siano stati soddisfatti. In tutti gli altri casi, il \ac{PEP} nega l’accesso.
  \item[permit-biased]
  Questo algoritmo è il duale di deny-biased.
\end{description}
Questi algoritmi evidenziano il fatto che le \emph{obligation} non solo influenzano il processo di autorizzazione, ma anche l'enforcement.
Si nota infine che gli errori causati da obbligazioni opzionali sono ignorati.


\section{Esempio di politica in FACPL}
\label{sec:Esempio_FACPL}
In questo semplice esempio di politica in \ac{FACPL}, due utenti possono interagire con il sistema attraverso delle richieste.
John può scrivere sulla risorsa "file.txt", ma non vi è specificata nessuna regola per la lettura. Invece Tom può
leggere il "file.txt", ma non può scriverci.\par
\lstinputlisting[language = FACPL, caption = {Esempio di politica in FACPL}\label{lst:facpl_es_Pol}]{./code/cap3/esempioFacpl}
Nel PAS si indicano gli algoritmi usati e le richieste da includere.
In questo esempio entrambi gli utenti richiedono sia l'azione di "WRITE" che l'azione di "READ".
Qui di seguito si presenta la struttura delle \emph{request}.
\lstinputlisting[language = FACPL, caption = {Esempio di richieste in FACPL}\label{lst:facpl_es_Req}]{./code/cap3/esempioRichiesteFacpl}
L'output dopo l'esecuzione sarà il seguente:
\begin{verbatim}
  Request: Request1

    Authorization Decision: PERMIT
    emph{obligation}: PERMIT M log_permit([John])

  Request: Request2

    Authorization Decision: NOT_APPLICABLE
    obligations:

  Request: Request3

    Authorization Decision: PERMIT
    obligations: PERMIT M log_permit([Tom])

  Request: Request4

    Authorization Decision: DENY
    obligations: DENY M log_deny([Tom])
\end{verbatim}
La prima richiesta di scrittura da parte di John viene chiaramente accettata, la sua richiesta di lettura però ha
come risultato \texttt{NOT APPLICABLE} in quanto non ci sono regole che possono essere usate per dare una valida risposta.
La richiesta di lettuta di Tom invece viene accettata, mentre la sua ultima richiesta di scrittura riceve un deny
perché è bloccata alla \emph{rule} writeRuleT.\par
Da questo esempio si può vedere la semplicità con cui si possono scrivere le policy e le richieste.
Le stesse politiche scritte in \ac{XACML} oltre a risultare prolisse a confronto, sono difficili da comprendere o analizzare.\par
Infine si fa notare che in questa versione, le richieste in ingresso sono completamente indipendenti tra loro e
la parte dello Usage Control descritta in Sezione \ref{sec:Usage_Control}, in cui si enfatizza l'uso di controlli continuativi e
la gestione di contesti mutabili, non è ancora implementata. Nel capitolo successivo si mostra la nuova struttura per
per uno sviluppo basato proprio sullo Usage Control.

\myChapter{Usage Control in FACPL}
\label{chap:Usage Control in FACPL}
In questo capitolo si descrive un'estensione del linguaggio \ac{FACPL} basata sul modello \ac{UCON}.
Più precisamente lo sviluppo è stato diviso in due parti. In primis la struttura è stata modificata di modo
che si possa creare delle politiche che si basino sul comportamento passato. In seguito il linguaggio
è stato ulteriormente esteso per migliorare la gestione dell'utilizzo continuativo di risorse .\par
In questo capitolo e nel succesivo si espone il secondo sviluppo, per la prima parte dell'estensione si rimanda
alla tesi del mio collega Federico Schipani, con cui ho lavorato per l'implementazione dell'intera nuova struttura. \par
Il linguaggio è stato esteso in modo tale da ottimizzare la gestione di un insieme di richieste sulla stessa risorsa.
Per far questo è stata necessaria una modifica sul processo di valutazione~(Sezione \ref{sec:Il processo di valutazione Facpl})
e l'inserimento di nuovi componenti~(Sezione~\ref{sec:SintassiNew}).
\section{Il processo di valutazione}
\label{sec:Il processo di valutazione Facpl}
Per memorizzare il comportamento passato, è stato essenziale l'aggiunta di un altro elemento nel processo di valutazione.
Come si vede in Figura \ref{fig:img/evaluationProcessStatus}, la nuova componente è Status.\par
\MyFigure{img/evaluationProcessStatus}{Processo di valutazione Status}{1}
Oltre agli \emph{attribute names} il Context Handler avrà anche la possibilità di ricavare gli \emph{status attributes}.
Il \ac{PDP} quindi potrà richiedere entrambi i tipi di attributi per elaborare la sua risposta. La decisione poi passerà al PEP
come avveniva precedentemente.\par
Inoltre, sono state aggiunte un tipo di obbligazioni che possono modificare lo Status. Il \ac{PEP} quindi può far eseguire questo tipo di
azioni per cambiare gli attributi.\par
Riprendendo il primo esempio di~\ref{subs:Accessi in lettura e scrittura di file}, si riporta \emph{Se uno degli amministratori
sta scrivendo su un file, nessuno può leggere o apportare modifiche nello stesso momento}. Questa regola è realizzabile
solo con l'utilizzo di uno status attribute (e.g.\ un booleano isWriting) che viene modificato da un'azione del PEP dopo che
una richiesta di scrittura è stata accettata. Il \ac{PDP} nelle successive richieste di lettura o di scritturà farà un controllo
sull'attributo e negherà di conseguenza l'accesso.\par
\MyFigure{img/evaluationProcessPEP}{Processo di valutazione PEP-Check}{1}
Il passo successivo consiste nell'eliminazione della ridondanza dei controlli da parte del \ac{PDP}.\par
Quando in \ref{subs:Accessi in lettura e scrittura di file}, si scrive:
\emph{Letture ripetute, su uno stesso insieme di file, vengono gestite in modo efficiente}
si fa riferimento al nuovo metodo per il controllo di richieste uniformi su una categoria di risorse e
in figura \ref{fig:img/evaluationProcessPEP} si mostra il processo di valutazione ridotto.\par
Si può definire un certo tipo di richieste che verrà interamente gestito dal \ac{PEP} dopo una prima, e unica,
valutazione del \ac{PDP}. Se un sistema deve gestire numerose richieste di lettura, e non ci sono elementi che bloccano
un accesso, è possibile accettare una richiesta esaminando soltanto la tipologia dell'azione e l'appartenenza di una risorsa ad
un insieme predefinito.\par
Prendiamo per esempio un sistema che gestisce dei file dentro delle cartelle. Supponiamo che ci sia una cartella
"Share" dove tutti gli utenti possono leggere solo se non ci sono utenti che scrivono.
Il sistema può gestire richieste di lettura ripetute, facendo un unico controllo per la verifica dell'assenza di
utenti che stanno scrivendo, alla prima domanda di accesso. Poi le richieste successive potranno essere
gestite dal \ac{PEP} con solo due controlli, uno sul tipo di azione, in questo caso una lettura, l'altro sull'appartenza
alla cartella "Share". Se i due controlli non sono validi si ritorna alla valutazione completa con il \ac{PDP}.
\section{Sintassi}
\label{sec:SintassiNew}
La modifica nella struttura ha reso necessario apportare dei cambiamenti anche nella grammatica.
La nuova sintassi è riportata in Tabella \ref{tab:facpl_new_syntax}.\par
Al \ac{PAS} è stato aggiunto \status della forma:
$$(status: Attribute^+)^?$$
quindi lo \status è un elemento opzionale e può essere formato da almeno un \emph{ Attribute}.

\emph{ Attribute} invece è della forma:
$$(Type\ Identifier (= Value)^?)$$
quindi è formato da un tipo, che può essere \emph{int}, \emph{float}, \emph{boolean} o \emph{date}, da una stringa
identificatrice e da un valore.\par
Le \emph{PepAction} sono state modificate in modo tale da eseguire operazioni sugli attributi dello stato
e sono della forma: \emph{nomeAzione}(\emph{Attribute},type)
per esempio l'azione di somma di interi è \emph{add}(\emph{Attribute},int).\par
Agli \emph{Attribute Names} è stata aggiunta la produzione \textit{Status/Identifier} per identificare i
termini dello Status nelle \emph{Expressions}.\par
La nuova produzione per le \emph{Obligations} della forma:
$$[Effect\ env:Expr\  status:Expr \ expiration:Value^?]$$
è utilizzata per la nuova gestione delle richieste da parte del \ac{PEP}. La nuova tipologia
non ha pepAction, in quanto il \ac{PEP} può modificare valori dello \status solo se il processo di valutazione è quello
completo, cioè quello che include anche i controlli eseguiti dal \ac{PDP}. Inoltre deve esserci un insieme di \emph{Expression}
affinchè il \ac{PEP} possa fare le proprie verifiche e può esserci un valore opzionale che indica il tipo della scadenza.
Da ora in poi il nuovo tipo di obbligazione sarà denominato \emph{Obligation Check.}\\
Infine \emph{expiration:$Value^?$} della forma:
$$Value ::= Int\ |\ Date$$
indicherà la scadenza delle \emph{Obligation Check}. Int determina il numero massimo di richieste, mentre Date
il limite temporale. Se CheckType viene omesso la \eobligation non ha scadenza.
\begin{table}
\footnotesize

\caption{Sintassi ausiliaria per le risposte}
\hrule
$
\begin{array}{@{\ }r@{\ \ \ \ }r@{\ }r@{\ \ }l@{\ }}

&&&\\[-.2cm]
{\textbf{PDP \ Responses }} &
\mathit{PDPResponse} & ::= & \langle \,\mathit{Decision} \ \ \ \mathit{FObligation}^* \rangle
\\[.2cm]
{\textbf{Decisions}} &
\mathit{Decision} & ::= & \permit \Sep \deny \Sep \notApp \Sep \indet
\\[.2cm]
{\textbf{Fulfilled obligations}} &
\mathit{FObligation} & ::= &  [ \, \mathit{ObType} \ \ \obl{\mathit{Value}} \, ]\\[.1cm]
&& & \Sep [ \, \x{env:} \, \mathit{Expr}\ \  \x{status:}\mathit{Expr}\, ]

\end{array}
$\\
\hrule
\label{tab:facpl_new_context_syntax}
\end{table}
{\label{Sintassi ausiliaria di FACPL Aggiornata}}
\begin{table}
\centering
\scriptsize
\caption{Sintassi di $FACPL_{PB}$}
\hrule
$
\begin{array}{@{\,}r@{\ \ }r@{\ }r@{\ \ }l@{\ }}

&&&\\[-.2cm]
{\textbf{Policy Authorisation Systems}} &
\mathit{PAS} & ::= & ( \,  \x{pep:} \, \mathit{EnfAlg}\ \ \x{pdp:}\, \mathit{PDP} \, \  \x{status:}\, \mathit{
[Attribute]})
\\[.2cm]
{\textbf{Attribute}} &
\mathit{Attribute}
& ::= & (\mathit{Type} \ \mathit{Identifier} \ (= Value )^{?})
\\[.2cm]
{\textbf{Type}} &
\mathit{Type}
& ::= & \x{int} \ | \ \x{boolean} \ | \ \x{date} \ | \ \x{double}
\\[.2cm]
{\textbf{Enforcement algorithms}} &
\mathit{EnfAlg}
& ::= & \based \Sep \denyBiased \Sep \permitBiased
\\[.2cm]
{\textbf{Policy Decision Points}} &
\mathit{PDP} & ::= & \pdpPol{\algNT\ }{\x{policies:} \, \mathit{Policy}^{+}}
\\[.2cm]
{\textbf{Combining algorithms}} &
\algNT & ::= & \permitOver \Sep \denyOver \Sep \denyUnless \Sep \permitUnless \\
&& \mid &
\firstApp \Sep \onlyOneApp \Sep \weakCon \Sep \strongCon
\\[.2cm]
{\textbf{Policies}} &
\mathit{Policy} & ::= &
\ruleOpt{\mathit{Effect}\ \ \x{target:} \, Expr\ \ \x{obl:} \, \mathit{Obligation}^{*} \, } \\
&& \mid &
\{ \algNT\ \ \x{target:} \, Expr\ \
\x{policies:} \, \mathit{Policy}^{+}  \ \ \x{obl:} \, \mathit{Obligation}^{*} \, \}
\\[.2cm]
{\textbf{Effects}} &
\mathit{Effect} & ::= & \permit \Sep \deny
\\[.2cm]
{\textbf{Obligations}} &
\mathit{Obligation} & ::= & [ \, \mathit{Effect} \ \ \mathit{ObType} \ \ \obl{Expr} \, ]\\
&& & \Sep [ \, \mathit{Effect} \ \ \x{env:} \, \mathit{Expr}\ \  \x{status:}\mathit{Expr} \ \ \x{exp:}\mathit{Value}^?\, ]
\\[.2cm]
{\textbf{PepAction}} & \mathit{PepAction} & ::= & \, \x{add(\mathit{Attribute}, int)} \Sep \x{flag(\mathit{Attribute}, boolean)} \\
&& & \Sep \x{sumDate(\mathit{Attribute}, date)} \Sep \x{div(\mathit{Attribute}, int)} \\
&& & \Sep \x{add(\mathit{Attribute}, float)} \ \Sep \x{mul(\mathit{Attribute}, float)} \\
&& & \Sep \x{mul(\mathit{Attribute}, int)} \ \Sep \x{div(\mathit{Attribute}, float)} \\
&& & \Sep \x{sub(\mathit{Attribute}, int)} \ \Sep \x{sub(\mathit{Attribute}, float)} \\
&& & \Sep \x{sumString(\mathit{Attribute}, string)} \\
&& & \Sep \x{setValue(\mathit{Attribute}, string)}\\
&& & \Sep \x{setDate(\mathit{Attribute}, date)}
\\[.2cm]
{\textbf{Obligation Types}} &
\mathit{ObType} & ::= & M \Sep O
\\[.4cm]
\textbf{Expressions}&
\mathit{Expr} & ::= &
\mathit{Name} \Sep \mathit{Value}  \\
& & \mid &\x{and(\mathit{Expr}, \mathit{Expr})} \Sep \x{or(\mathit{Expr}, \mathit{Expr})} \Sep \x{not(\mathit{Expr})} \\
& & \mid &
 \x{equal(\mathit{Expr},\mathit{Expr})}  \Sep \x{in}(\mathit{Expr}, \mathit{Expr}) \\
& & \mid & \x{greater}\textrm{-}\x{than(\mathit{Expr},\mathit{Expr})} \Sep \x{add(\mathit{Expr} ,\mathit{Expr} )}\\
& & \mid & \x{subtract(\mathit{Expr} ,\mathit{Expr} )} \Sep \x{divide(\mathit{Expr} ,\mathit{Expr} )}\\
& & \mid & \x{multiply(\mathit{Expr} ,\mathit{Expr} )}  \Sep \x{less}\textrm{-}\x{than(\mathit{Expr}, \mathit{Expr})}\\
\\[.2cm]
%
\textbf{Attribute Names} &
\mathit{Name} & ::= & \mathit{Identifier}/\mathit{Identifier} \ | \ \x{status}/\mathit{Identifier}\\[.2cm]
%
\textbf{Literal Values} &
\mathit{Value} & ::= & \x{true} \mid \x{false} \mid \mathit{Double} \mid \mathit{String} \mid \mathit{Date}
\\[.4cm]
{\textbf{Requests}} &
\mathit{Request} & ::= & {\attribute{\mathit{Name}}{\mathit{Value}}}^{+}
\\[.1cm]
\end{array}
$
\hrule
\label{tab:facpl_new_syntax}
\end{table}
{\label{Sintassi di FACPL Aggiornata}}

La sintassi delle risposte è rimasta quasi interamente invariata. Nella tabella \ref{tab:facpl_new_context_syntax}
si può vedere che le decisioni non sono cambiate, mentre c'è una nuova produzione per le \emph{Fulfilled obligations}.

La nuova forma indica le \emph{Fulfilled obligations} di tipo Check. Queste hanno solo due \emph{Expression}, una per
l'\emph{Environment} e una per lo \emph{Status}. I due termini esprimono i controlli che deve effettuare il \ac{PEP}
prendendo i valori dal \emph{Context Handler}. Le nuove \emph{Fulfilled obligations} sono prive di pepAction in quanto,
come enunciato prima, il \ac{PEP} non può eseguire azioni, ma solo verifiche sugli attributi.\par
\vspace{3mm}
Si propone adesso un esempio di una Policy che utilizza il nuovo tipo di obbligazione.
Nella Policy si vuole specificare che un utente di nome \emph{Charlie} può effetture una lettura solo se nessuno sta scrivendo
e vuole utilizzare la risorsa \emph{contabilita.xlsx}.
Se la richiesta viene accettata si deve cambiare l'attributo \emph{isReading} a true, oltre a ciò al susseguirsi
di un altra richiesta di lettura il sistema cambiarà il tipo di processo valutativo e il \ac{PEP} effettuerà i controlli
solo sul tipo di azione e sul nome della risorsa.\par
\lstinputlisting[language = FACPL, caption = {Esempio per la sintassi}\label{lst:esempio_sintassi}]{code/cap4/charlieRead}
\lstinputlisting[language = FACPL, caption = {Esempio PAS}\label{lst:esempio_sintassi}]{code/cap4/charliePAS}
\vspace{1em}
Si fa notare che con questa porzione di codice:
\lstinputlisting[firstline = 8, lastline = 8,language = FACPL, caption = {Esempio per la sintassi}\label{lst:esempio_sintassi}]{code/cap4/charlieRead}
\vspace{1em}
si esprime la nuova struttura per la creazione delle \emph{Obligation Check} ed è questo il costrutto
in cui si indicano i controlli che il \ac{PEP} deve effettuare. Come si vede, a differenza delle altre
obbligations, non è specificato il tipo M o O (mandatory o optional) in quanto tutte le obligation check
devono essere eseguite in ogni caso.
\section{Semantica}
\label{sub:Semantica_Ext}
La modifica della sintassi ha determinato anche variazioni nella semantica descritta in \ref{sec:sem_fpl}.\par
Il \ac{PDP} fa uso degli \emph{Status Attributes} per la valutazione di una richiesta in input.
Questi attributi sono modificabili tramite alcune azioni dette \emph{PepAction}, come l'operazione di somma \emph{add(Attribute,int)}
oppure l'assegnamento di una data con \emph{setDate(Attribute,date)}.
Sono definiti nel \ac{PAS} come descritto nell'esempio precedente con questa struttura:
\lstinputlisting[firstline = 8, lastline = 8,language = FACPL, caption = {Esempio status attributes}\label{lst:esempio_sintassi}]{code/cap4/charliePAS}
\par Il \ac{PEP} oltre a verificare che le PepAction siano eseguite, con l'aggiunta delle \emph{Obbligation Check}, ha il
controllo nella gestione delle richieste continuative. Se nella policy è presente una obbligazione che fa parte del nuovo tipo,
il \ac{PEP} ha il compito di controllare che le \emph{Expression} contenute nell'\emph{Obligation Check} siano vere.\par
Nell'esempio precedente l'obbligazione è definita con:
\lstinputlisting[firstline = 8, lastline = 8,language = FACPL, caption = {Esempio Obligation Check}\label{lst:esempio_sintassi}]{code/cap4/charlieRead}
il \ac{PEP} deve verificare quindi che l'azione richiesta sia una "read" e che la risorsia sia "contabilita.xlsx".
Se il controllo dà esito positivo la richiesta è subito accettata, se la richiesta non passa una delle verifiche si continuerà
il processo di valutazione ripartendo dal \ac{PDP} e svolgendo tutti i passi.\par
Le \emph{Obligation Check} possono essere permanenti, limitate temporalmente oppure limitate sul numero di richieste
accettate.\par
\begin{description}[labelindent=5pt,style=multiline,leftmargin=3.5cm]
  \item[Permanenti]Se sono del primo tipo e la verifica sulle \emph{Expression} è sempre vera, il \ac{PEP} continuerà a gestire le
  richieste senza il controllo del \ac{PDP}. La gestione cambia solo quando una richiesta non è conforme a quelle accettate.\par
  \item[Scadenza sul tempo]Se sono del secondo tipo, anche se la verifica sulle \emph{Expression} è vera, quando il tempo limite è stato superato,
  la gestione delle richieste cambia e si riprende il processo valutativo completo.
  \item[Scadenza sul numero di richieste]Il terzo tipo di \emph{Obligation Check} è simile al secondo, con l'unica differenza che il limite non è temporale,
  ma sul numero di richieste elaborate.
\end{description}

\section{Esempi di Usage Control}
\label{sec:Esempi}
In questa sessione si mostrano gli esempio descritti in \ref{subs:es_UC} utilizzando le nuove funzionalità.
\subsection{Accessi in lettura e scrittura di file}
\label{sub:RW_Code_Sec}
Nel sistema ci sono due utenti: Alice e Bob. Alice fa parte del gruppo degli Administator,
mentre Bob appartine ad un gruppo di default non specificato. Solo gli amministratori possono scrivere,
però tutti possono leggere file precedentemente specificati.

Se Alice fa una richiesta di scrittura, ha l'obbligo di cambiare il booleano \emph{isWriting} in \texttt{true}.
Se finisce di scrivere invece deve fare una richiesta di \emph{stopWrite} e assegnare \texttt{false} all'attributo.
Ovviamente se non sta scrivendo non può richiedere la fine dell'azione quindi \emph{stopWrite} è accettabile solo quando
isWriting è \texttt{true}.

Tutti gli utenti possono leggere i file solo se nessuno sta scrivendo su uno di questi e
se i file appartengono ad un determinato insieme. Nella policy che riguarda le letture si specifica che i
file debbano essere "thesis.tex" oppure "facpl.pdf", inoltre viene utilizzato il nuovo tipo di obbligation.

Il \ac{PAS} è differente dalla versione usata in \ref{sec:Esempio_FACPL} solo nella definizione dello status
in cui si inizializza isWriting a \texttt{false}.
\clearpage
\lstinputlisting[language = FACPL, caption = {Policy per esempio lettura e scrittura }\label{lst:esempio_RWP}]{code/cap4/checkReadWritePolicy}
Di seguito si mostrano 8 richieste per descrivere tutti i possibili scenari.
Le prime tre \emph{request} richiedono una lettura sia da parte di Bob che di Alice. Queste tre richieste sono simili tra loro,
differiscono solo per i file richiesti e tutte saranno accette perché rispettano Read\_Policy. La \emph{request1} anche
se essenzialmente uguale alle altre due passa per un processo di valutazione diverso. Infatti solo la prima richiesta
sarà valutata dal \ac{PDP} e dal \ac{PEP} mentre le altre solo da quest'ultimo. Nella prima si controlla che l'azione sia "read",
il file sia nell'insieme \{"thesis.tex","facpl.pdf"\} e che il Boolean isWriting sia falso. Nella seconda e nella terza
invece il \ac{PEP} verifica solo che l'azione sia una lettura e che il file richiesto sia valido.
I controlli del \ac{PEP} si esprimono in questa riga delle policy:
\lstinputlisting[firstline = 24, lastline = 24,language = FACPL, caption = {Policy per esempio lettura e scrittura }\label{lst:esempio_RWP}]{code/cap4/checkReadWritePolicy}
\lstinputlisting[firstline = 1, lastline = 15,language = FACPL, caption = {Richieste per esempio lettura e scrittura }\label{lst:esempio_RWR}]{code/cap4/checkReadWriteRequests}
\begin{verbatim}
  ---------------------------------------------------
  REQUEST N: 1
  REQUEST: read_request Bob
   Decision=
   PERMIT
  time 53ms
  ---------------------------------------------------
  REQUEST N: 2
  REQUEST: read_request Bob
   Decision=
   PERMIT
  time 6ms
  ---------------------------------------------------
  REQUEST N: 3
  REQUEST: read_request Alice
   Decision=
   PERMIT
  time 9ms
\end{verbatim}
La quarta richiesta interrompe il processo valutativo ridotto in quanto l'azione è una "write" e si ritorna alla valutazione
completa includendo i controlli del \ac{PDP}. La \emph{request5} e la \emph{request6} non vengono accettate perché l'amministratore
sta scrivendo sul file e solo successivamente alla richiesta di "stopWrite" di Alice, Bob può leggere il documento.
L'ultima richiesta è una lettura, ed essendo successiva ad un'altra read, si adotterà nuovamente il processo di valutazione ridotto.
\lstinputlisting[firstline = 16, language = FACPL, caption = {Richieste per esempio lettura e scrittura }\label{lst:esempio_RWR}]{code/cap4/checkReadWriteRequests}
\begin{verbatim}
  ---------------------------------------------------
  REQUEST N: 4
  REQUEST: write_request Alice
   Decision=
   PERMIT
  time 47ms
  ---------------------------------------------------
  REQUEST N: 5
  REQUEST: read_request Bob
   Decision=
   DENY
  time 25ms
  ---------------------------------------------------
  REQUEST N: 6
  REQUEST: stop_write_request Alice
   Decision=
   PERMIT
  time 44ms
  ---------------------------------------------------
  REQUEST N: 7
  REQUEST: read_request Bob
   Decision=
   PERMIT
  time 32ms
  ---------------------------------------------------
  REQUEST N: 8
  REQUEST: read_request Alice
   Decision=
   PERMIT
  time 3ms
\end{verbatim}
\subsection{Servizio di streaming}
\label{sub:Stream_Code_Sec}
L'esempio seguente è simile al primo, ma in questo caso si fa uso anche delle \emph{Obligations Check} limitate dal tempo.
Ci sono due utenti. Alice è un utente \emph{Premium} e può ascoltare le canzoni senza limiti una volta fatto il login. Bob
è un utente standard e dopo un certo lasso di tempo non può più fare richieste di ascolto e dovrà sentire la pubblicità prima
di poter ascoltare di nuovo un brano.

Le "listen" dei due tipi di utenti usano obligations check differenti. Gli ascolti degli utenti premium sono associati
ad una obbligazione persistente. Possono quindi richiedere i brani senza essere limitati per tutta la durata del login.
Gli utenti standard invece possono ascoltare le canzoni per quindici minuti di fila come viene definito in questa riga:
\lstinputlisting[firstline = 44, lastline=44,language = FACPL, caption = {Policy per esempio servizio streaming }\label{lst:esempio_SP}]{code/cap4/checkStreamingPolicy}
Tutti e due gli ascolti sono gestiti interamente dal \ac{PEP} a partire dalla seconda richiesta consecutiva, come avveniva per
le letture nell'esempio precedente.
\lstinputlisting[language = FACPL, caption = {Policy per esempio servizio streaming }\label{lst:esempio_SP}]{code/cap4/checkStreamingPolicy}
Le prime due \emph{request} non sono accettate perché non è possibile richiedere gli ascolti prima di fare un login.
Dopo le richieste di login, dove si usa un attributo password per validare l'accesso, i due utenti possono ascoltare i brani.

Si suppone che dopo tre richieste di ascolto, Bob alla \emph{request8} abbia superato il limite di tempo di quindici minuti.
Tutte le sue possibili richieste di ascolto sono rifiutate finchè non richiede di sentire la pubblicità. Dopo una richiesta
di "listenCommercials" Bob ha di nuovo altri quindici minuti di ascolto possibili.
\lstinputlisting[language = FACPL, caption = {Richieste per esempio servizio streaming }\label{lst:esempio_SR}]{code/cap4/checkStreamingRequests}
\begin{verbatim}
  ---------------------------------------------------
  REQUEST N: 1
  REQUEST: listen alice
  PEP Decision=
   DENY
  ---------------------------------------------------
  REQUEST N: 2
  REQUEST: listen bob
  PEP Decision=
   DENY
  ---------------------------------------------------
  REQUEST N: 3
  REQUEST: login alice
  PEP Decision=
   PERMIT
  ---------------------------------------------------
  REQUEST N: 4
  REQUEST: login bob
  PEP Decision=
   PERMIT
  ---------------------------------------------------
  REQUEST N: 5
  REQUEST: listen alice
  PEP Decision=
   PERMIT
  ---------------------------------------------------
  REQUEST N: 6
  REQUEST: listen bob
  PEP Decision=
   PERMIT
  ---------------------------------------------------
  REQUEST N: 7
  REQUEST: listen bob
  PEP Decision=
   PERMIT
  ---------------------------------------------------
  REQUEST N: 8
  REQUEST: listen bob
  PEP Decision=
   DENY
  ---------------------------------------------------
  REQUEST N: 9
  REQUEST: commercials bob
  PEP Decision=
   PERMIT
  ---------------------------------------------------
  REQUEST N: 10
  REQUEST: listen bob
  PEP Decision=
   PERMIT
\end{verbatim}

\myChapter{Estensione della libreria FACPL}
\label{chap:Estensione della libreria FACPL}
La librearia di supporto a \ac{FACPL} è basata su Java. Per implementare le nuove funzioni sono state aggiunte varie classi e sono state modificate
alcune parti già esistenti per adattare il comportamento ai nuovi componenti e alle nuove strutture.\par
Nelle sezioni successive si descriveranno gli aspetti più rilevanti del codice. L'intera libreria si può trovare su
GitHub all'indirizzo
\begin{center}
  \url{https://github.com/andreamargheri/FACPL}
\end{center}
nel branch Usage Control.
Lo sviluppo dell'estensione si può dividere in due passi principali.
La prima parte~(Serzione \ref{sec:PEPCheck}) consiste nel creare una classe \ac{PEP} che può gestire le richieste e adattare il processo di valutazione.
La seconda parte~(Sezioni \ref{sec:Le Obligation Check} e \ref{sec:FulfilledObligationCheck})
invece comprende l'implementazione di una gerarchia di \emph{Obbligations} che possono essere
sfruttate dal \ac{PEP} per la nuova gestione delle \emph{request}. \par
Nella Sezione~\ref{sec:Plugin Eclipse} viene mostrato come è stato esteso il plugin di FACPL
per supportare la nuova estensione.
Nella Sezione~\ref{sec:Esempi in Java} si descrivono le classi Java associate ai file \ac{FACPL}
dell'esempio in Sezione \ref{sub:Stream_Code_Sec}.
Infine, nella Sezione \ref{sec:Valutazione delle prestazioni} si mostrano i miglioramenti
prestazionali utilizzando l'esempio sviluppato in Sezione \ref{sub:RW_Code_Sec}.
\section{Classe PEPCheck}
\label{sec:PEPCheck}
Il normale flusso del processo valutativo passa prima dal \ac{PEP}. Questo riceve la richiesta e poi la manda al
\ac{PDP} affinché possa valutarla. Infine la risposta del \ac{PDP} è passata al \ac{PEP} che controlla le obbligazioni e
elabora la decisione conclusiva.
Il processo di valutazione ridotto avviene invece interamente nel \ac{PEP} che usa solo il \emph{Context Handler}
per ricavare gli attributi dell'ambiente e dello stato.

Avendo bisogno di entrambe le modalità di valutazione, una per avere la decisione completa e l'altra per eliminare la
ridondanza di verifiche da parte del \ac{PDP} su richieste ripetute,
la soluzione più appropriata è stata quella di creare un monitor
in cui si potesse valutare l'esigenza d'uso di uno o dell'altro flusso valutativo.

La classe PEPCheck è il monitor che esegue i controlli per determinare la gestione più adatta al contesto.
Questa estende la classe \ac{PEP} e contiente come campo privato il \ac{PDP} e la sua autorizzazione.
Nella figura \ref{fig:img/classDiagramPep.png} si mostra il grafico \ac{UML} della classe PEPCheck e i relativi elementi
associati.\par
\MyF{img/classDiagramPep.png}{Diagramma PEPCheck}{1}{15cm}
Nel Codice \ref{lst:PEPCostruttore} si mostrano i campi della classe PEPCheck e il costruttore.
\javaFirstlineC{PEPCheck.java}{Costruttore PEPCheck}{22}{34}{PEPCostruttore}
L'elemento che contraddistingue di più la nuova classe dal precedente \ac{PEP} è la lista di
\emph{AbstractFulfilledObligationCheck} che è usata nel metodo \emph{doAuthorisation}
per prendere la decisione sul flusso valutativo. I casi da controllare sono tre:
\begin{description}

  \item[\normalfont\bfseries{\MakeUppercase{N}on ci sono \emph{Obligation Check} in lista}]
  {
  \hfill \\
  Non si apporta nessun cambiamento al normale flusso del processo valutativo.
  Si invia la richiesta al \ac{PDP} che poi la passa la sua decisione al PEP.
  Il codice successivo mostra la condizione che rende possibile il paradigma.
  \javaLastline{PEPCheck.java}{53}{59}{PEPNoCheck}
  }

  \item[\normalfont\bfseries{\MakeUppercase{N}ella lista c'è almeno una \emph{Obligation Check}}]
  {
  \hfill \\
  Il sistema valuta le \emph{Expression} associate alle obbligazioni con il metodo
  \emph{doPEPCheck} (vedi \ref{sec:Il processo di valutazione Facpl}
  per la struttura). Nel caso in cui la decisione sia \texttt{PERMIT} oppure \texttt{DENY},
  il risultato viene subito restituito senza i controlli del \ac{PDP}.
  \javaLastline{PEPCheck.java}{60}{69}{PEPCheck}
  }

  \item[\normalfont\bfseries{\MakeUppercase{I}l controllo sulla lista restituisce \texttt{error/\MakeUppercase{NOT\_APPLICABLE}}}]
  {
  \hfill \\
  Se la decisione non ricade nei casi precedenti,
  si è verificato un errore nella valutazione delle espressioni. Se ci sono problemi nei controlli, il PEPCheck, prima di
  inviare la risposta, esegue una valutazione completa facendo \emph{doAuthorisation} sul \ac{PDP} e \emph{doEnforcement}.
  \javaFirstlineN{PEPCheck.java}{67}{79}{PEPError}{28}
  }

\end{description}

Il \emph{doEnforcement} nel codice \ref{lst:PEPError} del PEPCheck richiama il metodo implementato in \ac{PEP}
nella linea \texttt{super.doEnforcement(authPDP, cxtReq);} come si mostra di seguito alla linea 48.
\javaLastline{PEPCheck.java}{82}{93}{PEPenforce}
Se la decisione non è \texttt{PERMIT} viene restituita subito la risposta negativa sull'autorizzazione e si conclude la valutazione.
Se invece la decisione è positiva si esegue un controllo sulle \emph{FulfilledObligation} e
nel caso in cui queste siano FulfilledObligationNumCheck oppure FulfilledObligationTimeCheck,
sono aggiunte alla lista di AbstractFulfilledObligationNumCheck di PEPCheck (linee 58-78). Solo dopo l'inserimento delle obbligazioni
si restituisce la decisione ricavata (linea 81).
\javaLastlineC{PEPCheck.java}{Aggiunta delle FulfilledObligationNumCheck}{95}{124}{PEPAddCheck}

La lista modificata nel Codice~\ref{lst:PEPAddCheck} è utilizzata nel metodo \emph{doPEPCheck} (nel Codice \ref{lst:PEPdoPEP})
per eseguire la valutazione di ogni singola obligation check senza passare dal \ac{PDP}.
Questo metodo controlla che i valori siano corretti confrontandoli con gli attributi del contesto. Se non ci
sono errori e le espressioni hanno superato i controlli, si ritiene la richiesta verificata
utilizzando il processo di valutazione ridotto. Nelle linee 90 e 98 si ricava il risultato delle
singole obbligazioni con il metodo \emph{getObligationResult} e si inserisce in lista \emph{decisionList},
alla linea 106 invece si utilizza l'algoritmo
di combinazione delle decisioni per valutare l'intera lista dei risultati.
\javaLastlineC{PEPCheck.java}{Valutazione FulfilledObligationCheck}{131}{156}{PEPdoPEP}

Nella Sezione successiva \ref{sec:Le Obligation Check} si descrivono le Obligation utilizzate
nel metodo \emph{doPEPCheck} Codice \ref{lst:PEPdoPEP}.

\section{Classe Obligation Check}
\label{sec:Le Obligation Check}
La libreria è stata estesa con le \emph{Obbligations Check} perché c'era bisogno di obbligazioni che avessero solo
\emph{Expressions} da verificare senza avere nella struttura le pepAction.
La \emph{Obbligation Check} è una classe che estende la preesistente \emph{AbstractObligation}.
In Figura \ref{fig:img/classDiagramObl.png} si mostra il diagramma \ac{UML} della classe ObligationCheck.
Nel Codice \ref{lst:OBLConstr} si mostra che la classe ha come campi due alberi di espressioni booleane, che possono
a loro volta contenere altre \emph{expressions}, due classi per indicare la scadenza e una variabile che determina
il tipo di obligation check.
\MyF{img/classDiagramObl.png}{Diagramma ObligationCheck}{1}{12cm}
\javaFirstlineC{ObligationCheck.java}{Costruttore e campi}{12}{31}{OBLConstr}
La scadenza, come già detto in Sezione \ref{sub:SemanticaNew}, può essere di tre tipi: persistente,
limitata temporalmente oppure limitata dal numero di richieste gestite. I tipi sono indicati con "P" per le obbligation Persistenti,
"T" per le Temporali e "N" per quelle Numeriche come si mostra nell'\texttt{enum} in Codice \ref{lst:OBLType}.
\javaFirstlineC{CheckObligationType.java}{Enum delle Obligation}{1}{5}{OBLType}
Per ogni obbligazione ci sono quattro costruttori in modo tale da poter creare l'oggetto per ogni combinazione di \emph{Expressions}
in input. In Sezione \ref{lst:OBLConstrT} si mostra un costruttore per le Obligation Temporali che prende in input due ExpressionBooleanTree,
la scadenza determinata da \emph{FacplDate expiration} e l'effetto (che può essere \texttt{PERMIT} oppure \texttt{DENY}).
\javaFirstlineC{ObligationCheck.java}{Costruttore Obligation Temporale}{127}{132}{OBLConstrT}
Il metodo \emph{init} nel Codice \ref{lst:OBLInit}, presente in tutti i costruttori, inizializza i campi della classe a seconda del tipo di scadenza data
in input. Se la scadenza non è presente si utilizza il metodo init per creare le obbligazioni di tipo Persistente.
\javaLastlineC{ObligationCheck.java}{Inizializzazione dei campi della classe}{59}{72}{OBLInit}

In \emph{createObligation} nel Codice \ref{lst:OBLFul} si utilizzano i campi assegnati dai costruttori
e da init per la creazione delle \emph{FulfilledObligationNumCheck}
che sono descritte più approfonditamente in Sezione \ref{sec:FulfilledObligationCheck}.
\javaLastlineC{ObligationCheck.java}{Creazione di FulfilledObligationNumCheck}{154}{170}{OBLFul}

\section{Classe Fulfilled Obligation Check}
\label{sec:FulfilledObligationCheck}
Le \emph{Fulfilled Obligation Check} sono una gerarchia di classi che è stata aggiunta
come estensione alla \emph{AbstractFulfilledObligation} già presente nella libreria.
Ogni tipo di Obligation check ha una Fulfilled Obligation Check associata. Quindi sono state create
tre classi differenti:
\begin{enumerate}
  \item FulfilledObligationNumCheck
  \item FulfilledObligationCheckPersistent
  \item FulfilledObligationTimeCheck
\end{enumerate}
Le tre \emph{Fulfilled Obligation Check} sono molto simili
tra loro e l'unica differenza è la loro gestione della scadenza.\par
In Figura \ref{fig:img/classDiagramFulfilled.png} si mostra il diagramma della gerarchia di classi.
\MyF{img/classDiagramFulfilled.png}{Diagramma FulfilledObligationCheck}{1.2}{22.5cm}
Quando si crea una Fulfilled Obligation con una scadenza determinata da un numero intero, si utilizza
la \emph{FulfilledObligationNumCheck}. Nel Codice \ref{lst:FULn} si mostra il costruttore di questo tipo di
FulfilledObligation nel caso in cui la prima espressione sia una \emph{ExpressionFunction}, la seconda
una \emph{ExpressionBooleanTree} e la scadenza sia, ovviamente, un intero.
\javaFirstlineC{FulfilledObligationCheck.java}{FulfilledObligationNumCheck N}{58}{66}{FULn}
\clearpage
Nel Codice \ref{lst:FULeV} e \ref{lst:FULeS} si mostra il funzionamento di \emph{GetObligationResult}.
Questo metodo ha il compito, sia di verificare il valore di
verità delle espressioni associate, sia di gestire il valore di \emph{expiration}.
\begin{description}
  \item[\normalfont\bfseries{\MakeUppercase{E}xpiration non scaduto}]
  {
  \hfill \\
  Se l'intero che identifica la scadenza è maggiore di zero,
  si possono valutare le \emph{Expressions} e richiamare il metodo \emph{subExpiration}(Codice \ref{lst:FULnSub})
  per aggiornare il valore.
  \javaLastlineC{FulfilledObligationCheck.java}{Condizione di expiration valido}{94}{100}{FULeV}
  }
  \item[\normalfont\bfseries{\MakeUppercase{E}xpiration scaduta}]
  {
  \hfill \\
  Se si raggiunge la scadenza (in questo caso quando expiration è zero), le expressions non sono valutate
  e si assegna direttamente il valore \texttt{ERROR}.
  \javaLastlineC{FulfilledObligationCheck.java}{Condizione di expiration scaduta}{101}{105}{FULeS}
  }
\end{description}
Nel Codice \ref{lst:FULeV} si richiama \emph{subExpiration}. Il metodo è utilizzato per aggiornare il
valore della scadenza. Come si vede nel Codice \ref{lst:FULnSub} se il valore è -1 non si modifica expiration
(si spiegherà il motivo in seguito nel paragrafo delle FulfilledObligationCheckPersistent).
Negli altri casi si sottrae una unità a \emph{expiration} e
se si raggiunge il valore zero, si assegna al campo \emph{hasExpired} \texttt{true}.
\javaLastlineC{FulfilledObligationCheck.java}{subExpiration}{147}{159}{FULnSub}

La Fulfilled persistente estende la classe FulfilledObligationNumCheck.
Il costruttore richiama la classe superiore e inizializza la scadenza numerica con
il valore speciale -1. Come si è mostrato in nel Codice \ref{lst:FULnSub}, se \emph{expiration= -1}
il metodo non modifica mai il valore. In questo modo in \emph{GetObligationResult} di
FulfilledObligationNumCheck non si raggiunge mai
il limite della scadenza, le funzioni non restituiscono mai il risultato \texttt{ExpressionValue.ERROR} associato
ad un obbligazione scaduta e si valuta sempre le \emph{Expressions}.
Nel Codice \ref{lst:FULp} si mostra il costruttore per FulfilledObligationCheck di tipo "P".
\javaFirstlineC{FulfilledObligationCheckPersistent.java}{FulfilledObligationCheck P}{8}{16}{FULp}

Infine se la scadenza è di tipo \emph{FACPLDate} allora si utilizzerà le FulfilledObligationTimeCheck.
Queste si comportano in modo analogo alle \emph{FulfilledObligationNumCheck}. La differenza che le contraddistingue è
l'utilizzo dei metodi \emph{before} e \emph{after} per i controlli sulla scadenza.
\javaFirstlineC{FulfilledObligationTimeCheck.java}{getObligationResult per le tipo T}{139}{156}{FULtGet}
Inoltre, al posto del metodo \emph{subExpiration} in \emph{GetObligationResult}, si richiama \emph{checkExpiration}.
La funzione è utilizzata in questo caso solo per richiamare il metodo di debug per mostrare
la differenza tra il tempo relativo alla determinata richiesta e la scadenza. In \ref{lst:FULtCheck} si
mostra come viene calcolato il divario tra la scadenza e il tempo.
\javaLastlineC{FulfilledObligationTimeCheck.java}{Time difference}{202}{205}{FULtCheck}
\section{Plugin Eclipse}
\label{sec:Plugin Eclipse}
\MyFigure{img/toolChain.png}{ToolChain di FACPL}{1}
La sintassi di \ac{FACPL} è definita in Xtext e viene usato Xtend per la traduzione da \ac{FACPL} in Java.
Xtext~\cite{xtext} è un framework open-source per lo sviluppo di linguaggi di programmazione.
Xtend~\cite{xtend} è un linguaggio di programmazione basato sintatticamente e semanticamente su Java, ma
ha una sintassi più concisa e funzionalità addizionali come: inferenza sui tipi, overload degli operatori e estensione dei metodi.
\MakeUppercase{è} compilato in codice Java, ha circa lo stesso sistema di tipi statico, ma offre inferenza e integra tutte le librerie Java esistenti.
Il plugin utilizza Xtext e Xtend per rendere Eclipse un ambiente di sviluppo per \ac{FACPL} che offre funzionalità
come l'autocompletamento e l'highlighting del codice. In Figura~\ref{fig:img/toolChain.png} si mostra come da un file \ac{FACPL}
si possa generare il codice Java o \ac{XACML} utilizzando le regole di traduzione scrtte in Xtend.\par
\MyFigure{img/XtextGram.png}{Produzione per le Obligation}{1}
\MyF{img/XtextObl.png}{Obligation in Xtend}{1}{9.2cm}
Inizialmente è stata estesa la grammatica aggiungendo le produzioni per i tre tipi di ObligationCheck.
Successivamente sono state modificate le regole di traduzione in Xtend in modo da tradurre il codice scritto in \ac{FACPL} in Java.
In Figura \ref{fig:img/XtextGram.png} si mostra la nuova produzione delle Obligation.
Dopo aver esteso la grammatica sono state aggiunte le regole di traduzione. In Figura \ref{fig:img/XtextObl.png} si può
vedere come verrà generata l'obbligazione nel file Java della politica in base al
tipo di obbligazione usata nel file \ac{FACPL}.\par
In Figura \ref{fig:img/XtextGen.png} si presenta il plugin. Questo offre funzionalità come l'autocompletamento, il
syntax highlighting, auto indentazione e il controllo sull'unicità dei nomi.
Inoltre si può generare codice Java o XACML basandosi su un file \ac{FACPL} usando il menù come si
vede in Figura \ref{fig:img/XtextGen.png}. Questo tool insieme ad una guida online si può trovare all'indirizzo:
\begin{center}
  \url{http://facpl.sourceforge.net}
\end{center}
\MyFigure{img/XtextGen.png}{Plugin di FACPL}{1}
\section{Esempi di Usage Control in Java}
\label{sec:Esempi in Java}
Si ripropongono adesso gli esempi descritti in \ref{sec:Esempi} mostrando le classi Java relative
ai file \ac{FACPL}.
\subsection{Servizio di streaming}
\label{sub:EsStreamingJava}
Le classi associate all'esempio di servizio streaming in Sezione \ref{sub:RW_Code_Sec}
sono StatusStreaming, le ContextRequest degli utenti Alice e Bob, PolicySet\_Streaming e MainFACPL.

La classe StatusStreaming contiene FacplStatus. Si fa notare che nel codice \ref{lst:statusStr},
lo status viene creato una sola volta e sarà unico per tutte le richieste. Se non è inizializzato
si creano gli attributi e si aggiungono al campo FacplStatus, mentre se è già inizializzato il metodo
\texttt{getStatus()} restituirà direttamente \texttt{status}.
\javaFirstlineC{checkStreamingExample/StatusStreaming.java}{Classe StatusStreaming}{9}{32}{statusStr}

Le richieste di ascolto e di login sono simili tra loro. Queste si differenziano in base agli attributi associati.
Nel Codice \ref{lst:requestStr} si mostra la richiesta di ascolto di Alice come esempio di \emph{request}.
Per prima cosa si creano due attributi (linee 9-10), uno per il nome dell'utente (Alice) e uno per il tipo di azione (listen).
Dopo la creazione, i due \emph{attributes} sono aggiunti alla \emph{request} (linee 12-13).
Infine si associa il contesto della richiesta allo Status con il metodo \texttt{setStatus()} (linea 21) che prende come attributo FacplStatus usando \texttt{getStatus()}
di StatusStreaming. \par
\javaFirstlineC{checkStreamingExample/ContextRequest_ListenPremiumAlice.java}{Request}{10}{33}{requestStr}

Il Codice \ref{lst:policyStr} è relativo alla Policy di ascolto dell'utente standard Bob.
Il target della policy è composto da due \emph{Expressions}. La prima è associata all'azione che deve essere una \texttt{listen},
la seconda invece si riferisce al nome dell'utente che deve essere \texttt{Bob} (linee 8-14).\par
Si fa notare la creazione di una FacplDate (new FacplDate("00:15:00")) di quindici minuti per assegnare la scadenza all'obligation check.
Nell'obbligazione si valuta inoltre che l'utente sia ancora \texttt{Bob} e che il booleano
\texttt{streamingBob} (inizializzato a \texttt{true} solo dopo il login) sia \texttt{true} (linee 21-30).
\javaFirstlineC{checkStreamingExample/PolicySet_Streaming.java}{Policy Set}{169}{200}{policyStr}
La regola da valutare associata al PolicySet è stata aggiunta con il metodo
\texttt{addPolicyElement(new Rule\_ListenBob())} (linea 19).
Nella \texttt{Rule\_ListenBob}, mostrata nel Codice \ref{lst:ruleStr}, si controlla che Bob abbia fatto il login e che
non debba ascoltare la pubblicità (linee 38-45). Dopo il controllo si modifica il booleano \texttt{commercialsBob} a true in modo tale
che, dopo la scadenza della ObligationCheck, Bob non possa fare un'altra richiesta di ascolto (linee 48-53).
\javaLastlineC{checkStreamingExample/PolicySet_Streaming.java}{\texttt{Rule\_ListenBob}}{202}{223}{ruleStr}

La classe MainFACPL nel Codice \ref{lst:mainStr} incorpora \ac{PDP}, PEPCheck e il PolicySet. Nel costruttore si inizializzano
il \ac{PDP} con un Combining Algorithm (in questo caso \texttt{PermitUnlessDenyGreedy}) e l'insieme di policy, il PEP con l'algoritmo di
Enforcement (\texttt{DENY\_BIASED}), Combining Algorithm (\texttt{DenyOverridesCheck}) e il \ac{PDP}.
\javaFirstlineC{checkStreamingExample/MainFACPL.java}{MainFACPL}{20}{35}{mainStr}
Nel metodo \texttt{main} si aggiungono le richieste nella lista \texttt{requests}, queste poi verrano
valutate nel ciclo proposto di seguito.
\javaLastline{checkStreamingExample/MainFACPL.java}{60}{72}{mainForStr}

\subsection{Accessi in lettura e scrittura di file}
\label{sub:EsAccesJava}
In questa sezione si mostrano le parti di codice Java più significative associate all'esempio
descritto in Sezione \ref{sub:RW_Code_Sec}.\par
\vspace{2mm}
I codici per le richieste e per la costruzione dello status in questo caso non saranno trattati
in quanto hanno la stessa struttura dei file precedenti e cambiano soltanto il nome degli attributi e il loro tipo.\par
Si presenta invece la policy sulle letture che utilizza l'obbligazione di tipo permanente. A differenza della
delle regole sull'ascolto di Bob mostrato nel Codice \ref{lst:policyStr}, l'\emph{ObligationCheck} della lettura nel Codice
\ref{lst:policyFileCheck} non ha nessuna scadenza, ma solo due \emph{Expressions} in cui si indica il controllo sull'
insieme dei file in cui è possibile fare la lettura e l'azione concessa(\texttt{read}) (linee 25-28). Con questo tipo di struttura
l'obbligazione non raggiunge mai la scadenza fintantoché gli attributi delle richieste passano la verifica delle due espressioni.
\clearpage
\javaFirstlineC{checkWriteExample/PolicySet_ReadWriteCheck.java}{Policy di lettura}{80}{107}{policyFileCheck}

\section{Valutazione delle prestazioni}
\label{sec:Valutazione delle prestazioni}
In questa sezione si mostrano le differenze tra i tempi computazionali relativi a un
insieme di policy che utilizza le \emph{ObligationCheck}
e uno che non le utilizza basandosi sull'esempio in Sezione \ref{sub:EsAccesJava}.\par
In Sezione \ref{sub:RW_Code_Sec} è stato mostrato il tempo di esecuzione per ogni singola richiesta e
si è notato che le richieste valutate utilizzando solo il \ac{PEP} ricevevano una risposta molto più velocemente rispetto
alle altre. Per questo motivo è stato ideato un semplice sistema di benchmark per presentare i miglioramenti a livello
prestazionale attraverso l'uso delle \emph{ObligationCheck} nel processo valutativo.\par
Nel Codice \ref{lst:policyFileNoCheck} si presenta la policy duale del Codice \ref{lst:policyFileCheck},
in questo caso le regole non fanno uso del nuovo tipo di obbligazione.
\javaLastlineC{checkWriteExample/PolicySet_ReadWrite.java}{Policy di lettura senza ObligationCheck}{83}{105}{policyFileNoCheck}

Nella classe MainFACPL si utilizzano i due insiemi di regole per controllare due liste di richieste differenti.
\javaFirstlineC{checkWriteExample/MainFACPL.java}{Main per i benchmarks}{53}{70}{mainBench}
Nella prima lista sono presenti cento richieste di lettura, nella seconda il numero totale di richieste è lo stesso,
ma ogni otto letture c'è una scrittura e una domanda di interruzione delle modifiche.\par
Nel Codice \ref{lst:mainBench} si mostra il caso in cui si utilizza il PolicySet senza ObligationCheck per la prima lista
(tutte letture). Tutti i tempi di computazione vengono scritti nel file "testNoCheck.txt" (linea 9).\par
\MyFigure{img/bench.png}{Computazione}{1}
Utilizzando il set di policy senza ObligationCheck il \ac{PDP} controlla sempre ogni richiesta. Se invece si usa l'insieme di regole
con la nuova obbligazione, nella prima lista di richieste, dopo il primo controllo, il \ac{PDP} non sarà più parte del processo
di valutazione, nella seconda lista invece il \ac{PDP} controllerà il 30\% delle richieste.
Ogni lista di richieste è stata verificata cento volte per entrambi i tipi di Policy per un totale di diecimila valutazioni.

In Figura \ref{fig:img/bench.png} le unità in verticale indicano i millisecondi impiegati per controllare cento richieste.
La linea blu e quella arancione indicano i tempi per la computazione utilizzando il PolicySet con le ObligationCheck
mentre le altre due indicano i tempi per l'insieme di regole che non utilizzano il nuovo paradigma.
Analizzando i tempi totali per valutare diecimila richieste, in Figura \ref{fig:img/benchTot.png} si può vedere che
per la prima lista (solo richieste di lettura) il sistema impiega tre volte meno tempo utilizzando le ObligationCheck, mentre per la
seconda lista c'è un miglioramento del 35\% circa.

Come era ragionevole attendersi, le misurazioni confermano che, se le richieste di scrittura sono minori rispetto a quelle di lettura, la nuova gestione
alleggerisce il processo valutativo.
\MyFigure{img/benchTot.png}{Tempo totale}{1}

\myChapter{Conclusioni}
\label{chap:Conclusioni}
In questa tesi è stata sviluppata un'estensione del linguaggio \ac{FACPL} per migliorare la gestione delle richieste di accesso.
Nella prima parte sono state descritte le strutture su cui si è basato lo sviluppo. Sono stati introdotti i primi modelli
dell'Access Control partendo dal più semplice \ac{ACL} arrivando fino a \ac{PBAC}. In seguito è stato mostrato lo Usage Control
descrivendo il nuovo modello UCON\textsubscript{ABC} creato da Jaehong Park e Ravi Sandhu.

Successivamente è stato esaminato il linguaggio \ac{FACPL} su cui si è basato implementazione del monitor a runtime.
\MakeUppercase{è} stata mostrata la sintassi, la semantica e il processo di valutazione per presentare i lati positivi
e soprattutto i limiti del linguaggio su cui si è lavorato per migliorare il controllo continuativo degli accessi.

Sono state aggiunte nuove produzioni alla sintassi ed è stata modificata la risposta dell'autorizzazione per poter
essere utilizzata nel nuovo processo valutativo, in cui si elimina i controlli superflui per un determinato tipo di richieste.
Alle modifiche apportate sulla sintassi sono seguite le implementazioni di nuove strutture nella libreria Java su cui si
basa \ac{FACPL}. Infine sono stati proposti esempi per presentare le operazioni possibili e i miglioramenti nell'uso delle
risorse computative.

\section{Sviluppi Futuri}
\label{sec:Sviluppi Futuri}
In questo documento sono stati esposti due esempi, in uno si usa il numero di richieste, nell'altro il tempo.
La libreria di \ac{FACPL} è semplice da estendere ed è facile pensare a nuovi possibili controlli su cui
le autorizzazioni possono basarsi. Si potrebbe associare le verifiche della richiesta alla posizione
geografica dell'utente, oppure alla congestione della rete aggiungendo semplicemente la gestione di un nuovo parametro nella
creazione delle obbligazioni e le rispettive funzioni che ne controllano il valore.

\section*{Acronimi}
\label{sec:Acronimi}
\begin{acronym}[cap2]
  \acro{ACL}{Access Control List}
  \acro{RBAC}{Role Based Access Control}
  \acro{ABAC}{Attribute Based Access Control}
  \acro{PBAC}{Policy Based Access Control}
  \acro{UCON}{Usage Control}
	\acro{PAS}{Policy Authorization System}
	\acro{PR}{Policy Repository}
	\acro{PDP}[PDP]{Policy Decision Point}
 	\acro{PEP}[PEP]{Policy Enforcement Point}
	\acro{OASIS}{Organization for the Advancement of Structured Information Standards}
 	\acro{XACML}[XACML]{eXtensible Access Control Markup Language}
 	\acro{XML}[XML]{eXtensible Markup Language}
 	\acro{FACPL}[FACPL]{Formal Access Control Policy Language}
	\acro{EBNF}{Extended Backus-Naur form}
	\acro{UML}{Unified Modeling Language}
\end{acronym}
\nocite{*}
\bibliographystyle{abbrv}
\bibliography{Chapters/bib}
%--------------------------------------------------------------
\end{document}
%--------------------------------------------------------------
